\documentclass[a4paper,12pt]{report}

% --- Packages pour la langue et l'encodage ---
\usepackage[utf8]{inputenc}
\usepackage[T1]{fontenc}
\usepackage[french]{babel}
\usepackage{lmodern}

% --- Packages pour la mise en page ---
\usepackage[a4paper,margin=2.5cm,top=3cm,bottom=3cm]{geometry}
\usepackage{setspace}
\onehalfspacing
\usepackage{parskip}
\usepackage{titlesec}
\usepackage{tocloft}
\usepackage{fancyhdr}
\usepackage{graphicx}
\usepackage{booktabs}
\usepackage{longtable}
\usepackage{tabularx}
\usepackage{array}
\usepackage{multirow}
\usepackage{colortbl}

% --- Packages pour la typographie et les couleurs ---
\usepackage{xcolor}
\usepackage{fontawesome5}
\usepackage{tcolorbox}

% --- Packages pour le code et la syntaxe ---
\usepackage{minted}
\usemintedstyle{friendly}
\setminted{
    fontsize=\small,
    frame=lines,
    framesep=2mm,
    baselinestretch=1.2,
    bgcolor=gray!5,
    linenos=false
}

% --- Définition des couleurs personnalisées ---
\definecolor{primary}{RGB}{41, 128, 185}
\definecolor{secondary}{RGB}{231, 76, 60}
\definecolor{accent}{RGB}{46, 204, 113}
\definecolor{dark}{RGB}{44, 62, 80}
\definecolor{light}{RGB}{236, 240, 241}
\definecolor{codebg}{RGB}{245, 245, 245}
\definecolor{jsonkey}{RGB}{156, 220, 254}
\definecolor{jsonstring}{RGB}{206, 145, 120}
\definecolor{jsonnumber}{RGB}{181, 206, 168}
\definecolor{jsonboolean}{RGB}{150, 150, 255}

% --- Configuration des en-têtes et pieds de page ---
\pagestyle{fancy}
\fancyhf{}
\fancyhead[L]{\textbf{Plateforme SaaS d'Intelligence Financière}}
\fancyhead[R]{\thepage}
\fancyfoot[C]{\textit{Rapport de Projet -- DATAERA}}
\renewcommand{\headrulewidth}{0.5pt}
\renewcommand{\footrulewidth}{0.5pt}

% --- Configuration des sections ---
\titleformat{\chapter}[display]
{\normalfont\huge\bfseries\color{primary}}
{\chaptertitlename\ \thechapter}{20pt}{\Huge}
\titleformat{\section}
{\normalfont\Large\bfseries\color{dark}}
{\thesection}{1em}{}
\titleformat{\subsection}
{\normalfont\large\bfseries\color{primary}}
{\thesubsection}{1em}{}

% --- Configuration de la table des matières ---
\renewcommand{\cftchapfont}{\bfseries\color{primary}}
\renewcommand{\cftsecfont}{\color{dark}}
\renewcommand{\cftsubsecfont}{\color{dark!70}}

% --- Commandes personnalisées ---
\newcommand{\api}[1]{\texttt{\color{primary}#1}}
\newcommand{\module}[1]{\textbf{\color{secondary}#1}}
\newcommand{\tech}[1]{\texttt{\color{accent}#1}}
\newcommand{\code}[1]{\texttt{\color{dark}#1}}

% --- Environnement pour les blocs JSON ---
\newtcolorbox{jsonbox}{
    colback=codebg,
    colframe=primary,
    boxrule=0.5pt,
    arc=2pt,
    left=6pt,
    right=6pt,
    top=6pt,
    bottom=6pt,
    title={\textbf{Spécification API}},
    coltitle=white,
    fonttitle=\bfseries
}

% --- Environnement pour les alertes ---
\newtcolorbox{alertbox}{
    colback=red!5,
    colframe=secondary,
    boxrule=0.5pt,
    arc=2pt,
    left=6pt,
    right=6pt,
    top=6pt,
    bottom=6pt,
    title={\textbf{\faExclamationTriangle\ Attention}},
    coltitle=white,
    fonttitle=\bfseries
}

% --- Environnement pour les informations ---
\newtcolorbox{infobox}{
    colback=blue!5,
    colframe=primary,
    boxrule=0.5pt,
    arc=2pt,
    left=6pt,
    right=6pt,
    top=6pt,
    bottom=6pt,
    title={\textbf{\faInfoCircle\ Information}},
    coltitle=white,
    fonttitle=\bfseries
}

% --- Environnement pour les bonnes pratiques ---
\newtcolorbox{bestpractice}{
    colback=green!5,
    colframe=accent,
    boxrule=0.5pt,
    arc=2pt,
    left=6pt,
    right=6pt,
    top=6pt,
    bottom=6pt,
    title={\textbf{\faCheckCircle\ Bonne Pratique}},
    coltitle=white,
    fonttitle=\bfseries
}

% --- Informations du document ---
\title{\textbf{\Huge Plateforme SaaS d'Intelligence Financière}\\
\large Conception et Analyse Architecturale}
\author{\textbf{Rapport de Stage Technique}\\
\textit{DATAERA -- Solutions Financières Intégrées}}
\date{\today}

\begin{document}

% --- Page de titre ---
\begin{titlepage}
    \centering
    \vspace*{2cm}
    \includegraphics[width=3cm]{example-image-a}\par\vspace{1cm}
    {\Huge \textbf{\color{primary}Plateforme SaaS d'Intelligence Financière}}\par\vspace{0.5cm}
    {\Large Conception et Analyse Architecturale}\par\vspace{1.5cm}
    {\large \textbf{Rapport de Projet Technique}}\par\vspace{0.5cm}
    {\textit{DATAERA -- Solutions Financières Intégrées}}\par\vspace{2cm}
    {\large \today}\par
    \vfill
\end{titlepage}

% --- Table des matières ---
\tableofcontents
\newpage

% ==============================================================================
% CHAPITRE 1 : INTRODUCTION
% ==============================================================================
\chapter{Introduction}

\section{Vision Globale de la Plateforme SaaS}

La transformation numérique du secteur financier a conduit à l'émergence de plateformes SaaS (Software as a Service) sophistiquées, capables d'automatiser et d'intelligence-iser les processus financiers complexes. Dans ce contexte, notre projet s'articule autour du développement d'une \textbf{Plateforme SaaS d'Intelligence Financière} globale, conçue pour répondre aux besoins croissants des entreprises en matière d'automatisation, de conformité et d'analyse financière.

\subsection{Contexte du Secteur Financier Numérique}

L'industrie financière connaît une mutation profonde sous l'impulsion de plusieurs facteurs convergents :

\begin{itemize}
    \item \textbf{Réglementations évolutives} : Les directives telles que PSD2 en Europe et les normes AML/CFT au niveau mondial imposent de nouvelles exigences de transparence et de sécurité.
    \item \textbf{Attentes client} : Les entreprises et les particuliers exigent des services financiers plus rapides, plus transparents et personnalisés.
    \item \textbf{Progrès technologiques} : L'IA, le Machine Learning et le Big Data permettent désormais d'automatiser des processus auparavant manuels.
    \item \textbf{Concurrence des Fintech} : Les nouveaux entrants disruptifs poussent les institutions traditionnelles à innover.
\end{itemize}

\subsection{Objectifs d'Automatisation}

La plateforme vise à accomplir plusieurs objectifs stratégiques :

\begin{itemize}
    \item \textbf{Centralisation des données financières} : Agrégation des données bancaires, comptables et transactionnelles dans un référentiel unique.
    \item \textbf{Automatisation intelligente} : Utilisation de l'IA et du Machine Learning pour automatiser les tâches répétitives et détecter les anomalies.
    \item \textbf{Conformité réglementaire} : Assurer le respect des normes AML/CFT (Anti-Money Laundering / Combating the Financing of Terrorism) et PSD2 (Payment Services Directive 2).
    \item \textbf{Détection proactive des risques} : Identification précoce des fraudes et des risques financiers grâce à des algorithmes avancés.
    \item \textbf{Reporting analytique} : Génération de rapports financiers détaillés et de tableaux de bord interactifs pour la prise de décision.
\end{itemize}

\subsection{Approche Architecturale}

L'architecture de la plateforme repose sur une approche de \textbf{microservices modulaires}, permettant une évolutivité horizontale et une maintenance simplifiée. Chaque module fonctionne de manière autonome tout en s'intégrant harmonieusement dans l'écosystème global via des API RESTful bien définies et des contrats d'interface OpenAPI.

\begin{infobox}
La plateforme adopte une architecture événementielle avec une couche de messagerie (Kafka/RabbitMQ) pour assurer la communication asynchrone entre les modules, garantissant ainsi résilience et scalabilité.
\end{infobox}

\subsection{Périmètre du Projet}

Ce rapport se concentre spécifiquement sur deux modules critiques de la plateforme :

\begin{itemize}
    \item \textbf{Module F2 -- Multi Banking Hub} : Responsable de l'agrégation et de la synchronisation des données bancaires via les protocoles PSD2 et EBICS.
    \item \textbf{Module F5 -- AI Fraud Detection} : Chargé de la détection proactive des fraudes financières par une approche hybride règles + Machine Learning.
\end{itemize}

Ces deux modules constituent le fondement technique de la plateforme, le premier assurant l'alimentation en données et le second garantissant la sécurité des transactions.

% ==============================================================================
% CHAPITRE 2 : CONTEXTE THÉORIQUE DES PROTOCOLES BANCAIRES
% ==============================================================================
\chapter{Contexte Théorique des Protocoles Bancaires}

\section{Fondamentaux de l'Open Banking}

\subsection{Historique et Évolution Réglementaire}

L'Open Banking représente une évolution majeure dans le secteur financier, initiée par la directive PSD2 (Payment Services Directive 2) adoptée par l'Union européenne en 2015. Cette réglementation vise à :

\begin{itemize}
    \item \textbf{Démocratiser l'accès aux données bancaires} : Les clients peuvent autoriser des tiers à accéder à leurs informations de compte.
    \item \textbf{Stimuler la concurrence} : Réduire le monopole des banques traditionnelles sur les données de paiement.
    \item \textbf{Améliorer l'expérience client} : Permettre l'émergence de services innovants et personnalisés.
    \item \textbf{Renforcer la sécurité} : Imposer l'authentification forte (SCA - Strong Customer Authentication).
\end{itemize}

\subsection{Architecture Technique de l'Open Banking}

L'architecture Open Banking repose sur plusieurs composants techniques standardisés :

\subsubsection{Spécifications Berlin Group}

Le Berlin Group a défini des spécifications techniques harmonisées pour l'Open Banking en Europe, incluant :

\begin{itemize}
    \item \textbf{XS2A Framework} : Spécifications pour l'accès aux informations de compte (XS2A - XS2A Interface).
    \item \textbf{Standardisation des endpoints} : Uniformisation des URLs et des formats de requêtes/réponses.
    \item \textbf{Gestion du consentement} : Mécanismes normalisés pour l'autorisation et la révocation des accès.
    \item \textbf{Sécurité} : Définition des mécanismes OAuth 2.0, mTLS, et signatures digitales.
\end{itemize}

\subsubsection{Flux d'Authentification OAuth 2.0}

Le processus d'authentification suit le flux Authorization Code avec PKCE (Proof Key for Code Exchange) :

\begin{enumerate}
    \item L'application redirige l'utilisateur vers la page de consentement de la banque.
    \item L'utilisateur s'authentifie et autorise l'accès à ses comptes.
    \item La banque redirige vers l'application avec un code d'autorisation.
    \item L'échange du code contre un access token et un refresh token.
    \item Utilisation de l'access token pour accéder aux APIs bancaires.
\end{enumerate}

\begin{bestpractice}
L'utilisation de PKCE est obligatoire pour les applications mobiles et recommandée pour les applications web afin de prévenir les attaques par interception de code d'autorisation.
\end{bestpractice}

\section{Fondamentaux du Protocole EBICS}

\subsection{Origine et Standardisation}

EBICS (Electronic Banking Internet Communication Standard) est un protocole de communication bancaire développé conjointement par l'Allemagne, la France et l'Autriche. Il est devenu le standard de facto pour les communications bancaires B2B en Europe.

\subsection{Architecture et Sécurité EBICS}

\subsubsection{Couche de Sécurité Multi-niveaux}

EBICS implémente une sécurité en profondeur à plusieurs niveaux :

\begin{itemize}
    \item \textbf{Transport Layer Security (TLS)} : Chiffrement de la communication réseau.
    \item \textbf{Signature XML} : Authentification et intégrité des messages échangés.
    \item \textbf{Chiffrement hybride} : Combinaison de chiffrement symétrique (pour les données) et asymétrique (pour les clés).
    \item \textbf{Gestion des clés} : Système de clés publiques/privées avec gestion du cycle de vie.
\end{itemize}

\subsubsection{Types de Clés EBICS}

Le protocole EBICS utilise plusieurs types de clés avec des rôles spécifiques :

\begin{table}[H]
\centering
\begin{tabularx}{\textwidth}{|X|X|X|}
\hline
\textbf{Type de Clé} & \textbf{Fonction} & \textbf{Utilisation} \\
\hline
INI & Initialisation & Échange initial des clés avec la banque \\
\hline
HIA & Authentification & Authentification des requêtes ultérieures \\
\hline
E002 & Signature & Signature des ordres bancaires \\
\hline
A006 & Chiffrement & Chiffrement des données sensibles \\
\hline
\end{tabularx}
\caption{Types de Clés EBICS et Leurs Fonctions}
\end{table}

\subsubsection{Ordres Bancaires EBICS}

Les ordres EBICS (Bank Orders) sont classés en plusieurs catégories :

\begin{itemize}
    \item \textbf{HKD (Hochkreditabruf Dialog)} : Téléchargement de relevés de compte.
    \item \textbf{HTD (Historischer Transaction Abruf Dialog)} : Téléchargement de l'historique des transactions.
    \item \textbf{HAA (Hochkreditabruf Abruf)} : Téléchargement de soldes et de positions.
    \item \textbf{CCT (Crédit Transfer)} : Virements de crédit.
    \item \textbf{STA (Status Abruf)} : Interrogation du statut d'ordres.
\end{itemize}

\begin{alertbox}
La gestion des clés EBICS nécessite une attention particulière à la sécurité. Les clés privées ne doivent jamais être stockées en clair et doivent être protégées par des HSM (Hardware Security Modules) ou des KMS (Key Management Services).
\end{alertbox}

\section{Comparaison PSD2 vs EBICS}

\subsection{Cas d'Usage Respectifs}

\begin{table}[H]
\centering
\begin{tabularx}{\textwidth}{|X|X|X|}
\hline
\textbf{Critère} & \textbf{PSD2 (Open Banking)} & \textbf{EBICS} \\
\hline
Public cible & B2C, Fintech, Applications mobiles & B2B, Entreprises, Trésorerie \\
\hline
Type d'accès & Consentement utilisateur & Agrément bancaire direct \\
\hline
Complexité & Moyenne (OAuth 2.0) & Élevée (signatures XML, gestion clés) \\
\hline
Volume de données & Transactions récentes & Historique complet, massivement \\
\hline
Latence & Temps réel à quasi-temps réel & Différé, par lots \\
\hline
Standardisation & Hautement standardisé & Standard national avec variations \\
\hline
\end{tabularx}
\caption{Comparaison PSD2 vs EBICS}
\end{table}

\subsection{Complémentarité des Protocoles}

Dans notre architecture, PSD2 et EBICS sont complémentaires :

\begin{itemize}
    \item \textbf{PSD2} pour les cas d'usage B2C et les applications nécessitant une expérience utilisateur fluide.
    \item \textbf{EBICS} pour les traitements de masse B2B et les besoins de réconciliation comptable.
    \item \textbf{Hybridation} : Utilisation combinée selon le profil client et les exigences opérationnelles.
\end{itemize}

\section{Considérations de Conformité Réglementaire}

\subsection{Normes AML/CFT}

La conformité aux normes Anti-Money Laundering (AML) et Combating the Financing of Terrorism (CFT) impose :

\begin{itemize}
    \item \textbf{KYC (Know Your Customer)} : Vérification de l'identité des clients.
    \item \textbf{Transaction Monitoring} : Surveillance continue des transactions suspectes.
    \item \textbf{Reporting} : Déclaration des activités suspectes aux autorités (STR - Suspicious Transaction Reports).
    \item \textbf{Sanctions Screening} : Vérification contre les listes de sanctions internationales.
\end{itemize}

\subsection{GDPR et Protection des Données}

Le règlement général sur la protection des données (RGPD) impose :

\begin{itemize}
    \item \textbf{Consentement explicite} : Autorisation claire pour l'utilisation des données.
    \item \textbf{Minimisation des données} : Collecte limitée aux données nécessaires.
    \item \textbf{Droit à l'oubli} : Possibilité de supprimer les données sur demande.
    \item \textbf{Portabilité} : Possibilité de transférer les données vers un autre service.
\end{itemize}

\begin{infobox}
La combinaison des exigences AML/CFT et GDPR crée parfois des tensions : les obligations de conservation pour la conformité AML peuvent entrer en conflit avec le droit à l'oubli du RGPD. Une politique de gestion du cycle de vie des données est essentielle.
\end{infobox}

% ==============================================================================
% CHAPITRE 3 : ARCHITECTURE GLOBALE ET SYNERGIE
% ==============================================================================
\chapter{Architecture Globale et Synergie}

\section{Vue d'Ensemble des Modules}

La plateforme se compose de sept modules fonctionnels interconnectés, chacun répondant à un besoin spécifique de l'écosystème financier. Mon travail s'est concentré sur les modules \module{F2 -- Multi Banking Hub} et \module{F5 -- AI Fraud Detection}, qui constituent le cœur du système d'agrégation de données et de sécurité.

\section{Modules des Collègues (Contexte Global)}

\subsection{Module F4 -- AI Accounting Intelligence}

Le \module{F4} assure la détection automatique des anomalies comptables au sein de l'entreprise. Ce module analyse les écritures comptables pour identifier :

\begin{itemize}
    \item Les doublons potentiels dans les saisies comptables
    \item Les incohérences entre les différents documents financiers
    \item Les écarts de régularisation automatiques
    \item Les anomalies de rapprochement bancaire
\end{itemize}

\subsection{Module F6 -- AI Financial Copilot}

Le \module{F6} propose un assistant conversationnel basé sur une architecture RAG (Retrieval-Augmented Generation). Ce copilot financier permet :

\begin{itemize}
    \item L'interrogation en langage naturel des données financières
    \item La génération d'explications contextuelles sur les transactions
    \item L'assistance à la décision pour les opérations financières
    \item L'automatisation des requêtes récurrentes via des templates
\end{itemize}

\subsection{Module F7 -- Financial Reporting \& Analytics}

Le \module{F7} assure la génération de rapports financiers et de tableaux de bord interactifs, incluant :

\begin{itemize}
    \item Des rapports de trésorerie en temps réel
    \item Des analyses de rentabilité par produit/service
    \item Des prévisions budgétaires basées sur l'historique
    \item Des visualisations interactives pour les directions financières
\end{itemize}

\section{Synergie et Intégration des Modules}

\subsection{Flux de Données Global}

La figure suivante illustre le flux de données entre les différents modules de la plateforme :

\begin{center}
\begin{tikzpicture}[node distance=2cm, auto, thick]
    \node[draw, rectangle, fill=primary!20, minimum width=2cm, minimum height=1cm] (f2) {F2 (Multi Banking)};
    \node[draw, rectangle, fill=secondary!20, minimum width=2cm, minimum height=1cm, right=of f2] (f5) {F5 (Fraud Detection)};
    \node[draw, rectangle, fill=accent!20, minimum width=2cm, minimum height=1cm, right=of f5] (f4) {F4 (Accounting)};
    \node[draw, rectangle, fill=yellow!20, minimum width=2cm, minimum height=1cm, below=of f5] (f6) {F6 (Copilot)};
    \node[draw, rectangle, fill=purple!20, minimum width=2cm, minimum height=1cm, below=of f4] (f7) {F7 (Reporting)};
    
    \draw[->] (f2) -- node[above] {\tiny Transactions} (f5);
    \draw[->] (f5) -- node[above] {\tiny Validées} (f4);
    \draw[->] (f4) -- (f7);
    \draw[->] (f5) -- (f6);
    \draw[->] (f4) -- (f6);
    \draw[->] (f7) -- (f6);
\end{tikzpicture}
\end{center}

\subsection{Rôle des Modules F2 et F5 dans l'Écosystème}

\subsubsection{Module F2 comme Source de Données}

Le \module{F2 -- Multi Banking Hub} joue un rôle fondamental en tant que \textbf{source primaire de données} pour l'ensemble de la plateforme. En centralisant les connexions bancaires via les protocoles Open Banking (PSD2) et EBICS, il alimente :

\begin{itemize}
    \item Le \module{F5} avec les transactions brutes pour l'analyse de fraude
    \item Le \module{F4} avec les relevés bancaires pour le rapprochement comptable
    \item Le \module{F7} avec les données de trésorerie pour les rapports financiers
    \item Le \module{F6} avec le contexte transactionnel pour les requêtes conversationnelles
\end{itemize}

\subsubsection{Module F5 comme Garde-Fou de Sécurité}

Le \module{F5 -- AI Fraud Detection} assure la \textbf{sécurité transversale} de la plateforme en :

\begin{itemize}
    \item Validant les transactions avant leur intégration comptable (F4)
    \item Fournissant des scores de risque pour les rapports financiers (F7)
    \item Enrichissant le contexte du copilot financier avec des alertes de fraude (F6)
    \item Alimentant le Multi Banking Hub avec des patterns de fraude connus
\end{itemize}

\begin{bestpractice}
L'intégration modulaire permet à chaque module de fonctionner de manière autonome tout en bénéficiant des données enrichies par les autres modules, créant ainsi un effet de synergie où la valeur globale dépasse la somme des valeurs individuelles.
\end{bestpractice}

% ==============================================================================
% CHAPITRE 4 : ARCHITECTURE MICROSERVICES ET PATTERNS AVANCÉS
% ==============================================================================
\chapter{Architecture Microservices et Patterns Avancés}

\section{Fondamentaux de l'Architecture Microservices}

\subsection{Principes de Conception}

L'architecture microservices repose sur plusieurs principes fondamentaux qui guident notre conception :

\begin{itemize}
    \item \textbf{Single Responsibility} : Chaque service est responsable d'une fonctionnalité métier spécifique.
    \item \textbf{Decoupling} : Services faiblement couplés avec des interfaces bien définies.
    \item \textbf{Autonomy} : Chaque service peut être développé, déployé et scalé indépendamment.
    \item \textbf{Fail Isolation} : L'échec d'un service ne doit pas entraîner l'arrêt du système complet.
    \item \textbf{Data Ownership} : Chaque service possède et gère sa propre base de données.
\end{itemize}

\subsection{Avantages et Défis}

\subsubsection{Avantages de l'Approche Microservices}

\begin{itemize}
    \item \textbf{Scalabilité granulaire} : Possibilité de scaler uniquement les services qui en ont besoin.
    \item \textbf{Déploiement continu} : Mises à jour fréquentes sans impacter l'ensemble du système.
    \item \textbf{Technological diversity} : Utilisation de technologies optimales pour chaque service.
    \item \textbf{Resilience} : Isolation des pannes et facilité de récupération.
    \item \textbf{Team autonomy} : Équipes petites et autonomes par service.
\end{itemize}

\subsubsection{Défis et Complexités}

\begin{itemize}
    \item \textbf{Complexité distribuée} : Gestion des communications entre services.
    \item \textbf{Data consistency} : Maintien de la cohérence des données distribuées.
    \item \textbf{Observability} : Monitoring et debugging de systèmes distribués.
    \item \textbf{Network latency} : Impact des communications réseau sur les performances.
    \item \textbf{Testing complexity} : Tests d'intégration plus complexes.
\end{itemize}

\section{Patterns de Communication Inter-Services}

\subsection{Communication Synchrone vs Asynchrone}

\subsubsection{RESTful APIs (Synchrone)}

Les APIs RESTful sont utilisées pour les communications synchrones où une réponse immédiate est attendue :

\begin{itemize}
    \item \textbf{Cas d'usage} : Requêtes utilisateur, validations en temps réel.
    \item \textbf{Avantages} : Simplicité, compatibilité HTTP standard, caching natif.
    \item \textbf{Inconvénients} : Couplage temporel, blocage possible, latence additive.
\end{itemize}

\subsubsection{Event-Driven Architecture (Asynchrone)}

L'architecture événementielle utilise des messages asynchrones pour la communication :

\begin{itemize}
    \item \textbf{Cas d'usage} : Traitements de fond, notifications, synchronisations.
    \item \textbf{Avantages} : Découplage temporel, résilience, scalabilité naturelle.
    \item \textbf{Inconvénients} : Complexité de gestion, ordering des messages, exactly-once semantics.
\end{itemize}

\begin{infobox}
Dans notre architecture, nous utilisons une approche hybride : REST pour les requêtes synchrones (ex: création de connexion) et Kafka pour les traitements asynchrones (ex: synchronisation bancaire).
\end{infobox}

\subsection{Patterns de Messagerie}

\subsubsection{Publish/Subscribe Pattern}

Le pattern Pub/Sub permet à plusieurs services de s'abonner à des événements spécifiques :

\begin{minted}{typescript}
// Publisher - Module F2
class BankSyncPublisher {
  constructor(private kafkaProducer: KafkaProducer) {}

  async publishTransactionSynced(event: TransactionSyncedEvent): Promise<void> {
    await this.kafkaProducer.send({
      topic: 'bank.transactions.synced',
      messages: [{
        key: event.connectionId,
        value: JSON.stringify(event),
        headers: {
          'event-type': 'TRANSACTION_SYNCED',
          'correlation-id': event.correlationId
        }
      }]
    });
  }
}

// Subscriber - Module F5
class FraudDetectionSubscriber {
  constructor(private kafkaConsumer: KafkaConsumer) {}

  async subscribe(): Promise<void> {
    await this.kafkaConsumer.subscribe({ topic: 'bank.transactions.synced' });
    
    await this.kafkaConsumer.run({
      eachMessage: async ({ message }) => {
        const event = JSON.parse(message.value.toString());
        await this.processTransactionForFraud(event);
      }
    });
  }
}
\end{minted}

\subsubsection{Message Queue Pattern}

Les files d'attente permettent de traiter les messages de manière séquentielle et garantissent l'ordre de traitement :

\begin{minted}{typescript}
class TransactionProcessor {
  private queue: PriorityQueue<Transaction> = new PriorityQueue();
  
  async enqueueTransaction(transaction: Transaction): Promise<void> {
    await this.queue.enqueue(transaction, transaction.priority);
  }
  
  async processQueue(): Promise<void> {
    while (!this.queue.isEmpty()) {
      const transaction = await this.queue.dequeue();
      await this.processTransaction(transaction);
    }
  }
}
\end{minted}

\section{Patterns de Résilience et Tolérance aux Pannes}

\subsection{Circuit Breaker Pattern}

Le Circuit Breaker empêche les cascades d'échecs en ouvrant un "circuit" lorsqu'un service échoue de manière répétée :

\begin{minted}{typescript}
enum CircuitState {
  CLOSED = 'CLOSED',    // Normal operation
  OPEN = 'OPEN',        // Failing, reject requests
  HALF_OPEN = 'HALF_OPEN' // Testing recovery
}

class CircuitBreaker {
  private state: CircuitState = CircuitState.CLOSED;
  private failureCount = 0;
  private lastFailureTime?: Date;
  private successCount = 0;

  constructor(
    private threshold: number = 5,
    private timeout: number = 60000 // 1 minute
  ) {}

  async execute<T>(operation: () => Promise<T>): Promise<T> {
    if (this.state === CircuitState.OPEN) {
      if (this.shouldAttemptReset()) {
        this.state = CircuitState.HALF_OPEN;
      } else {
        throw new Error('Circuit breaker is OPEN - service unavailable');
      }
    }

    try {
      const result = await operation();
      this.onSuccess();
      return result;
    } catch (error) {
      this.onFailure();
      throw error;
    }
  }

  private onSuccess(): void {
    this.failureCount = 0;
    if (this.state === CircuitState.HALF_OPEN) {
      this.successCount++;
      if (this.successCount >= 2) {
        this.state = CircuitState.CLOSED;
        this.successCount = 0;
      }
    }
  }

  private onFailure(): void {
    this.failureCount++;
    this.lastFailureTime = new Date();
    
    if (this.failureCount >= this.threshold) {
      this.state = CircuitState.OPEN;
    }
  }

  private shouldAttemptReset(): boolean {
    return this.lastFailureTime && 
      Date.now() - this.lastFailureTime.getTime() > this.timeout;
  }
}
\end{minted}

\subsection{Retry Pattern with Exponential Backoff}

Le retry avec backoff exponentiel permet de retenter les opérations défaillantes tout en évitant de surcharger le système :

\begin{minted}{typescript}
class RetryHandler {
  async executeWithRetry<T>(
    operation: () => Promise<T>,
    maxRetries: number = 3,
    baseDelay: number = 1000
  ): Promise<T> {
    let attempt = 0;
    
    while (attempt < maxRetries) {
      try {
        return await operation();
      } catch (error) {
        attempt++;
        
        if (attempt === maxRetries) {
          throw new Error(`Operation failed after ${maxRetries} attempts: ${error.message}`);
        }
        
        const delay = baseDelay * Math.pow(2, attempt);
        const jitter = Math.random() * 0.1 * delay; // Add jitter
        await this.sleep(delay + jitter);
      }
    }
    
    throw new Error('Max retries exceeded');
  }

  private sleep(ms: number): Promise<void> {
    return new Promise(resolve => setTimeout(resolve, ms));
  }
}
\end{minted}

\subsection{Bulkhead Pattern}

Le Bulkhead Pattern isole les ressources pour empêcher un composant défaillant d'affecter l'ensemble du système :

\begin{minted}{typescript}
class BulkheadExecutor {
  private semaphore: Semaphore;
  
  constructor(maxConcurrent: number) {
    this.semaphore = new Semaphore(maxConcurrent);
  }

  async execute<T>(operation: () => Promise<T>): Promise<T> {
    const permit = await this.semaphore.acquire();
    
    try {
      return await operation();
    } finally {
      this.semaphore.release(permit);
    }
  }
}

class Semaphore {
  private permits: number;
  private queue: Array<() => void> = [];

  constructor(permits: number) {
    this.permits = permits;
  }

  async acquire(): Promise<void> {
    if (this.permits > 0) {
      this.permits--;
      return Promise.resolve();
    }
    
    return new Promise(resolve => {
      this.queue.push(resolve);
    });
  }

  release(): void {
    this.permits++;
    if (this.queue.length > 0) {
      const next = this.queue.shift();
      this.permits--;
      next();
    }
  }
}
\end{minted}

\section{Patterns de Gestion des Données}

\subsection{Database per Service Pattern}

Chaque microservice possède sa propre base de données, garantissant l'autonomie et le découplage :

\begin{itemize}
    \item \textbf{Avantages} : Découplage complet, indépendance technologique, isolation des données.
    \item \textbf{Défis} : Transactions distribuées, cohérence éventuelle, requêtes jointes complexes.
\end{itemize}

\subsection{Saga Pattern pour Transactions Distribuées}

Le Saga Pattern gère les transactions distribuées en les décomposant en une séquence de transactions locales avec des mécanismes de compensation :

\begin{minted}{typescript}
interface SagaStep<T> {
  execute: (data: T) => Promise<any>;
  compensate: (data: T, result: any) => Promise<void>;
}

class SagaOrchestrator<T> {
  constructor(private steps: SagaStep<T>[]) {}

  async execute(initialData: T): Promise<void> {
    const results: any[] = [];
    
    for (let i = 0; i < this.steps.length; i++) {
      const step = this.steps[i];
      
      try {
        const result = await step.execute(initialData);
        results.push(result);
      } catch (error) {
        // Compensation en cas d'échec
        for (let j = i - 1; j >= 0; j--) {
          try {
            await this.steps[j].compensate(initialData, results[j]);
          } catch (compensationError) {
            console.error(`Compensation failed for step ${j}:`, compensationError);
          }
        }
        throw error;
      }
    }
  }
}

// Exemple : Création de connexion bancaire
const bankConnectionSaga = new SagaOrchestrator([
  {
    execute: async (data) => {
      // Étape 1 : Créer l'enregistrement de connexion
      return await createConnectionRecord(data);
    },
    compensate: async (data, result) => {
      // Annuler la création de connexion
      await deleteConnectionRecord(result.connectionId);
    }
  },
  {
    execute: async (data) => {
      // Étape 2 : Initialiser le consentement OAuth
      return await initiateOAuthConsent(data);
    },
    compensate: async (data, result) => {
      // Révoquer le consentement
      await revokeOAuthConsent(result.consentId);
    }
  },
  {
    execute: async (data) => {
      // Étape 3 : Configurer les webhooks
      return await configureWebhooks(data);
    },
    compensate: async (data, result) => {
      // Désactiver les webhooks
      await disableWebhooks(result.webhookId);
    }
  }
]);
\end{minted}

\begin{bestpractice}
Le Saga Pattern avec compensation est particulièrement adapté aux transactions longue durée et aux processus métier complexes dans le domaine bancaire où les transactions ACID traditionnelles ne sont pas applicables.
\end{bestpractice}

\section{Patterns de Scalabilité}

\subsection{Horizontal Pod Autoscaling (HPA)}

L'autoscaling horizontal permet d'ajouter automatiquement des instances en fonction de la charge :

\begin{minted}{yaml}
# Kubernetes HPA Configuration
apiVersion: autoscaling/v2
kind: HorizontalPodAutoscaler
metadata:
  name: bank-connection-hpa
spec:
  scaleTargetRef:
    apiVersion: apps/v1
    kind: Deployment
    name: bank-connection-service
  minReplicas: 2
  maxReplicas: 10
  metrics:
  - type: Resource
    resource:
      name: cpu
      target:
        type: Utilization
        averageUtilization: 70
  - type: Resource
    resource:
      name: memory
      target:
        type: Utilization
        averageUtilization: 80
\end{minted}

\subsection{Database Sharding}

Le sharding partitionne les données horizontalement sur plusieurs instances de base de données :

\begin{minted}{typescript}
class ShardRouter {
  private shards: Map<string, DatabaseConnection> = new Map();
  
  constructor(shardConfig: ShardConfig[]) {
    shardConfig.forEach(config => {
      this.shards.set(config.shardId, new DatabaseConnection(config.connectionString));
    });
  }

  getShard(key: string): DatabaseConnection {
    const shardId = this.calculateShardId(key);
    return this.shards.get(shardId);
  }

  private calculateShardId(key: string): string {
    // Hash-based sharding
    const hash = this.hashFunction(key);
    const shardIndex = hash % this.shards.size;
    return Array.from(this.shards.keys())[shardIndex];
  }

  private hashFunction(key: string): number {
    let hash = 0;
    for (let i = 0; i < key.length; i++) {
      hash = ((hash << 5) - hash) + key.charCodeAt(i);
      hash |= 0; // Convert to 32bit integer
    }
    return Math.abs(hash);
  }
}

% ==============================================================================
% CHAPITRE 5 : MEILLEURES PRATIQUES DE SÉCURITÉ FINANCIÈRE
% ==============================================================================
\chapter{Meilleures Pratiques de Sécurité Financière}

\section{Fondamentaux de la Sécurité Financière}

\subsection{Principes de Sécurité (CIA Triad)}

La sécurité financière repose sur trois principes fondamentaux :

\begin{itemize}
    \item \textbf{Confidentialité} : Protection des données sensibles contre l'accès non autorisé.
    \item \textbf{Intégrité} : Garantie que les données ne sont pas altérées de manière non autorisée.
    \item \textbf{Disponibilité} : Assurance que les systèmes et données sont accessibles lorsque nécessaires.
\end{itemize}

\subsection{Cadre Réglementaire et Normes}

\subsubsection{PCI DSS (Payment Card Industry Data Security Standard)}

Le standard PCI DSS impose des exigences de sécurité pour la gestion des données de cartes de paiement :

\begin{itemize}
    \item \textbf{Protection des données} : Chiffrement des données de carte en transit et au repos.
    \item \textbf{Gestion des vulnérabilités} : Maintien de systèmes sécurisés et patchés.
    \item \textbf{Contrôle d'accès} : Restriction de l'accès aux données par besoin de savoir.
    \item \textbf{Monitoring} : Surveillance et test réguliers des réseaux.
\end{itemize}

\subsubsection{SOC 2 (Service Organization Control 2)}

Le framework SOC 2 définit des critères pour la gestion des services cloud :

\begin{itemize}
    \item \textbf{Sécurité} : Protection contre l'accès non autorisé.
    \item \textbf{Disponibilité} : Performance et disponibilité du service.
    \item \textbf{Intégrité du traitement} : Exhaustivité, exactitude et ponctualité du traitement.
    \item \textbf{Confidentialité} : Protection des informations confidentielles.
    \item \textbf{Privacy} : Collecte, utilisation, conservation et suppression des données personnelles.
\end{itemize}

\section{Authentification et Autorisation}

\subsection{OAuth 2.0 et OpenID Connect}

\subsubsection{Architecture OAuth 2.0}

OAuth 2.0 est un framework d'autorisation qui permet aux applications d'obtenir un accès limité aux comptes utilisateur :

\begin{itemize}
    \item \textbf{Resource Owner} : L'utilisateur qui autorise l'accès.
    \item \textbf{Client} : L'application demandant l'accès.
    \item \textbf{Authorization Server} : Serveur qui émet les tokens d'accès.
    \item \textbf{Resource Server} : Serveur qui héberge les données protégées.
\end{itemize}

\subsubsection{Implémentation de la Validation JWT}

\begin{minted}{typescript}
import jwt from 'jsonwebtoken';
import jwksClient from 'jwks-rsa';

class JWTValidator {
  private client: any;

  constructor(jwksUri: string) {
    this.client = jwksClient({
      jwksUri: jwksUri,
      cache: true,
      rateLimit: true,
      jwksRequestsPerMinute: 10
    });
  }

  async validateToken(token: string): Promise<DecodedToken> {
    try {
      const decoded = jwt.decode(token, { complete: true });
      
      if (!decoded || !decoded.header) {
        throw new Error('Invalid token format');
      }

      const key = await this.getSigningKey(decoded.header.kid);
      
      const verified = jwt.verify(token, key, {
        algorithms: ['RS256'],
        issuer: process.env.JWT_ISSUER,
        audience: process.env.JWT_AUDIENCE,
        clockTolerance: 30 // Tolérance de 30 secondes pour le décalage d'horloge
      });

      return verified as DecodedToken;
      
    } catch (error) {
      throw new Error(`Token validation failed: ${error.message}`);
    }
  }

  private async getSigningKey(kid: string): Promise<string> {
    return new Promise((resolve, reject) => {
      this.client.getSigningKey(kid, (err: any, key: any) => {
        if (err) {
          reject(err);
        } else {
          resolve(key.publicKey || key.rsaPublicKey);
        }
      });
    });
  }
}
\end{minted}

\subsection{Role-Based Access Control (RBAC)}

\subsubsection{Modèle de Permissions}

\begin{minted}{typescript}
enum Permission {
  // Banking Hub Permissions
  READ_BANK_CONNECTIONS = 'read:bank_connections',
  CREATE_BANK_CONNECTIONS = 'create:bank_connections',
  DELETE_BANK_CONNECTIONS = 'delete:bank_connections',
  SYNC_BANK_DATA = 'sync:bank_data',
  
  // Fraud Detection Permissions
  READ_FRAUD_CASES = 'read:fraud_cases',
  RESOLVE_FRAUD_CASES = 'resolve:fraud_cases',
  CONFIGURE_FRAUD_RULES = 'configure:fraud_rules',
  VIEW_FRAUD_ANALYTICS = 'view:fraud_analytics',
  
  // Admin Permissions
  MANAGE_USERS = 'manage:users',
  MANAGE_ROLES = 'manage:roles',
  VIEW_AUDIT_LOGS = 'view:audit_logs'
}

enum Role {
  SUPER_ADMIN = 'super_admin',
  BANK_ADMIN = 'bank_admin',
  FRAUD_ANALYST = 'fraud_analyst',
  ACCOUNTANT = 'accountant',
  READ_ONLY = 'read_only'
}

const ROLE_PERMISSIONS: Record<Role, Permission[]> = {
  [Role.SUPER_ADMIN]: Object.values(Permission),
  [Role.BANK_ADMIN]: [
    Permission.READ_BANK_CONNECTIONS,
    Permission.CREATE_BANK_CONNECTIONS,
    Permission.DELETE_BANK_CONNECTIONS,
    Permission.SYNC_BANK_DATA
  ],
  [Role.FRAUD_ANALYST]: [
    Permission.READ_FRAUD_CASES,
    Permission.RESOLVE_FRAUD_CASES,
    Permission.VIEW_FRAUD_ANALYTICS
  ],
  [Role.ACCOUNTANT]: [
    Permission.READ_BANK_CONNECTIONS,
    Permission.SYNC_BANK_DATA
  ],
  [Role.READ_ONLY]: [
    Permission.READ_BANK_CONNECTIONS,
    Permission.READ_FRAUD_CASES
  ]
};

class AuthorizationManager {
  constructor(private requiredPermissions: Permission[]) {}

  authorize(user: any): boolean {
    if (!user.roles || !Array.isArray(user.roles)) {
      return false;
    }

    const userPermissions = new Set<Permission>();
    
    user.roles.forEach((role: Role) => {
      const rolePermissions = ROLE_PERMISSIONS[role] || [];
      rolePermissions.forEach(perm => userPermissions.add(perm));
    });

    return this.requiredPermissions.every(permission =>
      userPermissions.has(permission)
    );
  }

  hasPermission(user: any, permission: Permission): boolean {
    const manager = new AuthorizationManager([permission]);
    return manager.authorize(user);
  }
}
\end{minted}

\section{Protection des Données Sensibles}

\subsection{Chiffrement des Données}

\subsubsection{Chiffrement au Repos (Encryption at Rest)}

\begin{minted}{typescript}
import crypto from 'crypto';

class DataEncryption {
  private algorithm = 'aes-256-gcm';
  private keyLength = 32; // 256 bits
  private ivLength = 16;  // 128 bits
  private tagLength = 16;  // 128 bits

  constructor(private encryptionKey: Buffer) {
    if (encryptionKey.length !== this.keyLength) {
      throw new Error('Encryption key must be 32 bytes (256 bits)');
    }
  }

  encrypt(plaintext: string): { encrypted: string; iv: string; tag: string } {
    const iv = crypto.randomBytes(this.ivLength);
    const cipher = crypto.createCipheriv(
      this.algorithm,
      this.encryptionKey,
      iv
    );

    let encrypted = cipher.update(plaintext, 'utf8', 'hex');
    encrypted += cipher.final('hex');
    
    const tag = cipher.getAuthTag();

    return {
      encrypted,
      iv: iv.toString('hex'),
      tag: tag.toString('hex')
    };
  }

  decrypt(encrypted: string, iv: string, tag: string): string {
    const decipher = crypto.createDecipheriv(
      this.algorithm,
      this.encryptionKey,
      Buffer.from(iv, 'hex')
    );

    decipher.setAuthTag(Buffer.from(tag, 'hex'));

    let decrypted = decipher.update(encrypted, 'hex', 'utf8');
    decrypted += decipher.final('utf8');

    return decrypted;
  }
}

// Utilisation avec AWS KMS
class KMSEncryption {
  private encryptionKey: Buffer;

  constructor() {
    // Récupération de la clé depuis AWS KMS
    this.encryptionKey = this.getDecryptedKeyFromKMS();
  }

  private async getDecryptedKeyFromKMS(): Promise<Buffer> {
    // Implémentation AWS KMS
    const kms = new AWS.KMS();
    const response = await kms.decrypt({
      CiphertextBlob: Buffer.from(process.env.ENCRYPTED_KEY, 'base64')
    }).promise();
    
    return response.Plaintext;
  }
}
\end{minted}

\subsubsection{Chiffrement en Transit (Encryption in Transit)}

\begin{minted}{typescript}
import https from 'https';
import fs from 'fs';

class SecureHTTPServer {
  constructor(private port: number) {}

  start(): void {
    const options = {
      key: fs.readFileSync(process.env.SSL_KEY_PATH),
      cert: fs.readFileSync(process.env.SSL_CERT_PATH),
      ca: fs.readFileSync(process.env.SSL_CA_PATH),
      minVersion: 'TLSv1.2', // Minimum TLS version
      ciphers: [
        'ECDHE-ECDSA-AES256-GCM-SHA384',
        'ECDHE-RSA-AES256-GCM-SHA384',
        'ECDHE-ECDSA-CHACHA20-POLY1305',
        'ECDHE-RSA-CHACHA20-POLY1305'
      ].join(':'),
      honorCipherOrder: true
    };

    const server = https.createServer(options, this.requestHandler);
    
    server.listen(this.port, () => {
      console.log(`Secure server running on port ${this.port}`);
    });
  }

  private requestHandler(req: any, res: any): void {
    // Implementation of request handler
    res.writeHead(200);
    res.end('Secure response');
  }
}
\end{minted}

\subsection{Masquage des Données (Data Masking)}

\begin{minted}{typescript}
class DataMasker {
  static maskIBAN(iban: string): string {
    if (!iban || iban.length < 8) return iban;
    
    const visibleChars = 4;
    const maskedChars = iban.length - visibleChars * 2;
    
    return iban.substring(0, visibleChars) + 
           '*'.repeat(maskedChars) + 
           iban.substring(iban.length - visibleChars);
  }

  static maskEmail(email: string): string {
    if (!email || !email.includes('@')) return email;
    
    const [local, domain] = email.split('@');
    const maskedLocal = local.charAt(0) + '*'.repeat(local.length - 1);
    
    return `${maskedLocal}@${domain}`;
  }

  static maskPhoneNumber(phone: string): string {
    if (!phone || phone.length < 4) return phone;
    
    const visibleChars = 2;
    const maskedChars = phone.length - visibleChars * 2;
    
    return phone.substring(0, visibleChars) + 
           '*'.repeat(maskedChars) + 
           phone.substring(phone.length - visibleChars);
  }

  static maskCreditCard(cardNumber: string): string {
    if (!cardNumber || cardNumber.length < 4) return cardNumber;
    
    const visibleChars = 4;
    const maskedChars = cardNumber.length - visibleChars;
    
    return '*'.repeat(maskedChars) + cardNumber.substring(cardNumber.length - visibleChars);
  }
}

// Utilisation dans les logs
class SecureLogger {
  static logTransaction(transaction: any): void {
    const maskedTransaction = {
      ...transaction,
      accountNumber: DataMasker.maskIBAN(transaction.accountNumber),
      customerEmail: DataMasker.maskEmail(transaction.customerEmail),
      cardNumber: DataMasker.maskCreditCard(transaction.cardNumber)
    };
    
    console.log(JSON.stringify(maskedTransaction));
  }
}
\end{minted}

\begin{alertbox}
Le masquage des données est essentiel dans les logs et les rapports de debugging pour éviter l'exposition accidentelle d'informations sensibles.
\end{alertbox}

\section{Sécurité des API}

\subsection{Rate Limiting et Throttling}

\begin{minted}{typescript}
import Redis from 'ioredis';

class RateLimiter {
  constructor(private redis: Redis) {}

  async checkRateLimit(
    identifier: string,
    limit: number,
    window: number // en secondes
  ): Promise<{ allowed: boolean; remaining: number; resetTime: number }> {
    const key = `ratelimit:${identifier}`;
    const current = await this.redis.incr(key);
    
    if (current === 1) {
      await this.redis.expire(key, window);
    }

    const ttl = await this.redis.ttl(key);
    const resetTime = Date.now() + (ttl * 1000);

    return {
      allowed: current <= limit,
      remaining: Math.max(0, limit - current),
      resetTime
    };
  }

  async checkSlidingWindow(
    identifier: string,
    limit: number,
    window: number
  ): Promise<{ allowed: boolean; remaining: number }> {
    const now = Date.now();
    const windowStart = now - (window * 1000);
    const key = `sliding:${identifier}`;

    // Supprimer les entrées expirées
    await this.redis.zremrangebyscore(key, 0, windowStart);

    // Compter les requêtes dans la fenêtre
    const count = await this.redis.zcard(key);

    if (count >= limit) {
      return { allowed: false, remaining: 0 };
    }

    // Ajouter la requête actuelle
    await this.redis.zadd(key, now, `${now}-${Math.random()}`);
    await this.redis.expire(key, window);

    return { allowed: true, remaining: limit - count - 1 };
  }
}

// Middleware Express
export function rateLimitMiddleware(limiter: RateLimiter, limit: number, window: number) {
  return async (req: any, res: any, next: any) => {
    const identifier = req.ip || req.user?.id || 'anonymous';
    
    const result = await limiter.checkRateLimit(identifier, limit, window);
    
    res.setHeader('X-RateLimit-Limit', limit);
    res.setHeader('X-RateLimit-Remaining', result.remaining);
    res.setHeader('X-RateLimit-Reset', result.resetTime);

    if (!result.allowed) {
      return res.status(429).json({
        error: 'Too Many Requests',
        message: 'Rate limit exceeded',
        retryAfter: Math.ceil((result.resetTime - Date.now()) / 1000)
      });
    }

    next();
  };
}
\end{minted}

\subsection{Input Validation et Sanitization}

\begin{minted}{typescript}
import { z } from 'zod';

// Schémas de validation avec Zod
const bankConnectionSchema = z.object({
  bank_id: z.string().min(1).max(50),
  connection_type: z.enum(['PSD2', 'EBICS']),
  credentials: z.object({
    client_id: z.string().min(10).max(100),
    client_secret: z.string().min(20).max(200),
    redirect_uri: z.string().url()
  }),
  scopes: z.array(z.enum(['ais', 'pis'])).min(1),
  consent_validity_days: z.number().int().min(1).max(365)
});

const fraudScanSchema = z.object({
  transaction_ids: z.array(z.string().uuid()).min(1).max(100),
  scan_types: z.array(z.enum(['aml_rules', 'ml_model', 'graph_analysis'])).min(1),
  risk_threshold: z.number().min(0).max(1),
  include_explanations: z.boolean().optional()
});

class InputValidator {
  static validateBankConnection(data: any): { valid: boolean; errors?: string[]; data?: any } {
    try {
      const validated = bankConnectionSchema.parse(data);
      return { valid: true, data: validated };
    } catch (error) {
      return {
        valid: false,
        errors: error.errors.map((e: any) => `${e.path.join('.')}: ${e.message}`)
      };
    }
  }

  static validateFraudScan(data: any): { valid: boolean; errors?: string[]; data?: any } {
    try {
      const validated = fraudScanSchema.parse(data);
      return { valid: true, data: validated };
    } catch (error) {
      return {
        valid: false,
        errors: error.errors.map((e: any) => `${e.path.join('.')}: ${e.message}`)
      };
    }
  }

  static sanitizeString(input: string): string {
    return input
      .replace(/[<>]/g, '') // Remove potential HTML tags
      .trim()
      .substring(0, 1000); // Limit length
  }

  static sanitizeSQL(input: string): string {
    // Basic SQL injection prevention
    const dangerousPatterns = [
      /('|(")|(;)|(\b(SELECT|INSERT|UPDATE|DELETE|DROP|UNION|EXEC)\b)/gi
    ];
    
    let sanitized = input;
    dangerousPatterns.forEach(pattern => {
      sanitized = sanitized.replace(pattern, '');
    });
    
    return sanitized;
  }
}
\end{minted}

\section{Audit et Conformité}

\subsection{Journalisation des Activités (Audit Logging)}

\begin{minted}{typescript}
interface AuditEvent {
  timestamp: Date;
  actor: string;
  action: string;
  resource: string;
  outcome: 'SUCCESS' | 'FAILURE';
  details?: any;
  ipAddress?: string;
  userAgent?: string;
  correlationId?: string;
}

class AuditLogger {
  constructor(private kafkaProducer: KafkaProducer) {}

  async logEvent(event: AuditEvent): Promise<void> {
    const enrichedEvent = {
      ...event,
      timestamp: event.timestamp.toISOString(),
      service: 'financial-platform',
      environment: process.env.NODE_ENV
    };

    await this.kafkaProducer.send({
      topic: 'audit.events',
      messages: [{
        key: event.actor,
        value: JSON.stringify(enrichedEvent),
        headers: {
          'event-type': 'AUDIT',
          'correlation-id': event.correlationId
        }
      }]
    });
  }

  async logBankConnectionCreated(
    actor: string,
    connectionId: string,
    bankId: string,
    correlationId: string
  ): Promise<void> {
    await this.logEvent({
      timestamp: new Date(),
      actor,
      action: 'BANK_CONNECTION_CREATED',
      resource: `bank-connection:${connectionId}`,
      outcome: 'SUCCESS',
      details: { bankId },
      correlationId
    });
  }

  async logFraudCaseResolved(
    actor: string,
    caseId: string,
    resolution: string,
    correlationId: string
  ): Promise<void> {
    await this.logEvent({
      timestamp: new Date(),
      actor,
      action: 'FRAUD_CASE_RESOLVED',
      resource: `fraud-case:${caseId}`,
      outcome: 'SUCCESS',
      details: { resolution },
      correlationId
    });
  }
}
\end{minted}

\subsection{Monitoring de la Sécurité}

\begin{minted}{typescript}
class SecurityMonitor {
  private suspiciousActivities: Map<string, number> = new Map();

  async detectAnomalousActivity(userId: string, activity: string): Promise<boolean> {
    const key = `${userId}:${activity}`;
    const count = (this.suspiciousActivities.get(key) || 0) + 1;
    this.suspiciousActivities.set(key, count);

    // Règles de détection d'anomalies
    const thresholds = {
      'failed_login': 5,
      'failed_transaction': 10,
      'api_abuse': 20
    };

    const threshold = thresholds[activity] || 15;
    
    if (count >= threshold) {
      await this.triggerSecurityAlert(userId, activity, count);
      return true;
    }

    return false;
  }

  private async triggerSecurityAlert(
    userId: string,
    activity: string,
    count: number
  ): Promise<void> {
    const alert = {
      timestamp: new Date(),
      severity: 'HIGH',
      type: 'SECURITY_ANOMALY',
      userId,
      activity,
      count,
      message: `Suspicious activity detected: ${activity} repeated ${count} times`
    };

    // Envoyer l'alerte au système de monitoring
    await this.sendToSecurityOps(alert);
    
    // Optionnellement, bloquer temporairement l'utilisateur
    await this.temporaryBlockUser(userId);
  }

  private async sendToSecurityOps(alert: any): Promise<void> {
    // Implémentation de l'envoi d'alerte
  }

  private async temporaryBlockUser(userId: string): Promise<void> {
    // Implémentation du blocage temporaire
  }
}

% ==============================================================================
% CHAPITRE 6 : MACHINE LEARNING POUR LA DÉTECTION DE FRAUDE
% ==============================================================================
\chapter{Machine Learning pour la Détection de Fraude}

\section{Fondamentaux du Machine Learning Financier}

\subsection{Types de Fraude Financière}

La détection de fraude financière nécessite une compréhension approfondie des différents types de menaces :

\begin{itemize}
    \item \textbf{Identity Theft} : Utilisation frauduleuse d'identités volées pour ouvrir des comptes ou effectuer des transactions.
    \item \textbf{Account Takeover} : Accès non autorisé à des comptes existants.
    \item \textbf{Card Not Present (CNP) Fraud} : Transactions par carte sans présentation physique.
    \item \textbf{Money Laundering} : Blanchiment d'argent via des transactions complexes.
    \item \textbf{Phishing} : Obtention d'informations sensibles par tromperie.
    \item \textbf{Insider Fraud} : Fraude commise par des employés ou partenaires internes.
\end{itemize}

\subsection{Défis Spécifiques au Domaine Financier}

\subsubsection{Déséquilibre des Classes (Class Imbalance)}

Les données de fraude présentent un déséquilibre extrême : typiquement moins de 0.1\% des transactions sont frauduleuses.

\begin{minted}{python}
import numpy as np
from imblearn.over_sampling import SMOTE
from imblearn.under_sampling import RandomUnderSampler
from imblearn.pipeline import Pipeline

class ImbalanceHandler:
    def __init__(self, sampling_strategy: str = 'auto'):
        self.sampling_strategy = sampling_strategy
    
    def handle_imbalance(self, X, y):
        """Gestion du déséquilibre des classes"""
        if self.sampling_strategy == 'smote':
            # Oversampling de la classe minoritaire
            smote = SMOTE(sampling_strategy='auto', random_state=42)
            X_resampled, y_resampled = smote.fit_resample(X, y)
            
        elif self.sampling_strategy == 'hybrid':
            # Approche hybride : oversampling + undersampling
            over = SMOTE(sampling_strategy=0.1, random_state=42)
            under = RandomUnderSampler(sampling_strategy=0.5, random_state=42)
            
            pipeline = Pipeline([('over', over), ('under', under)])
            X_resampled, y_resampled = pipeline.fit_resample(X, y)
            
        else:
            # Utilisation de class weights
            return X, y
            
        return X_resampled, y_resampled
    
    def calculate_class_weights(self, y):
        """Calcul des poids de classes pour le training"""
        from sklearn.utils.class_weight import compute_class_weight
        classes = np.unique(y)
        weights = compute_class_weight('balanced', classes=classes, y=y)
        return dict(zip(classes, weights))
\end{minted}

\subsubsection{Concept Drift}

Les patterns de fraude évoluent constamment, nécessitant des modèles adaptatifs :

\begin{minted}{python}
class ConceptDriftDetector:
    def __init__(self, window_size: int = 1000, threshold: float = 0.1):
        self.window_size = window_size
        self.threshold = threshold
        self.reference_window = None
        self.current_window = []
    
    def add_sample(self, sample: dict, prediction: float):
        """Ajoute un échantillon à la fenêtre courante"""
        self.current_window.append({
            'sample': sample,
            'prediction': prediction,
            'timestamp': datetime.now()
        })
        
        if len(self.current_window) >= self.window_size:
            self.check_drift()
    
    def check_drift(self):
        """Détecte le concept drift en comparant les distributions"""
        if self.reference_window is None:
            self.reference_window = self.current_window.copy()
            self.current_window = []
            return False
        
        # Comparaison des distributions using KS test
        from scipy import stats
        
        ref_predictions = [item['prediction'] for item in self.reference_window]
        curr_predictions = [item['prediction'] for item in self.current_window]
        
        statistic, p_value = stats.ks_2samp(ref_predictions, curr_predictions)
        
        drift_detected = p_value < self.threshold
        
        if drift_detected:
            print(f"Concept drift detected! KS statistic: {statistic}, p-value: {p_value}")
            self.reference_window = self.current_window.copy()
        
        self.current_window = []
        return drift_detected
\end{minted}

\section{Ingénierie des Caractéristiques (Feature Engineering)}

\subsection{Caractéristiques Transactionnelles}

\begin{minted}{python}
import pandas as pd
import numpy as np
from datetime import datetime, timedelta

class TransactionFeatureEngineer:
    def __init__(self):
        self.customer_profiles = {}
        self.merchant_profiles = {}
    
    def extract_features(self, transaction: dict) -> dict:
        """Extrait les caractéristiques d'une transaction"""
        features = {}
        
        # Caractéristiques temporelles
        features['hour_of_day'] = transaction['timestamp'].hour
        features['day_of_week'] = transaction['timestamp'].weekday()
        features['is_weekend'] = 1 if features['day_of_week'] >= 5 else 0
        features['is_night'] = 1 if features['hour_of_day'] >= 22 or features['hour_of_day'] <= 6 else 0
        
        # Caractéristiques de montant
        features['amount'] = transaction['amount']
        features['amount_log'] = np.log1p(transaction['amount'])
        features['is_high_amount'] = 1 if transaction['amount'] > 1000 else 0
        
        # Caractéristiques de localisation
        features['country'] = transaction['country']
        features['is_foreign'] = 1 if transaction['country'] != 'FR' else 0
        features['is_high_risk_country'] = 1 if transaction['country'] in self.high_risk_countries else 0
        
        # Caractéristiques de device
        features['device_type'] = transaction['device_type']
        features['is_mobile'] = 1 if transaction['device_type'] == 'mobile' else 0
        features['is_new_device'] = 1 if transaction['device_id'] not in self.customer_profiles.get(transaction['customer_id'], {}).get('known_devices', set()) else 0
        
        return features
    
    def extract_behavioral_features(self, customer_id: str, transaction: dict, historical_transactions: list) -> dict:
        """Extrait les caractéristiques comportementales basées sur l'historique"""
        features = {}
        
        if not historical_transactions:
            return features
        
        # Statistiques de montant
        amounts = [t['amount'] for t in historical_transactions]
        features['avg_amount'] = np.mean(amounts)
        features['std_amount'] = np.std(amounts)
        features['median_amount'] = np.median(amounts)
        features['amount_deviation'] = (transaction['amount'] - features['avg_amount']) / (features['std_amount'] + 1e-6)
        
        # Fréquence des transactions
        features['transaction_frequency'] = len(historical_transactions) / 30  # par jour sur 30 jours
        features['time_since_last_transaction'] = (transaction['timestamp'] - historical_transactions[-1]['timestamp']).total_seconds() / 3600
        
        # Patterns de localisation
        countries = [t['country'] for t in historical_transactions]
        features['unique_countries'] = len(set(countries))
        features['country_diversity'] = features['unique_countries'] / len(historical_transactions)
        
        # Patterns de device
        devices = [t['device_id'] for t in historical_transactions]
        features['unique_devices'] = len(set(devices))
        features['device_diversity'] = features['unique_devices'] / len(historical_transactions)
        
        return features
    
    def extract_network_features(self, transaction: dict, graph_data: dict) -> dict:
        """Extrait les caractéristiques de réseau/graphe"""
        features = {}
        
        # Caractéristiques de centralité
        features['customer_degree'] = graph_data.get('customer_degree', {}).get(transaction['customer_id'], 0)
        features['merchant_degree'] = graph_data.get('merchant_degree', {}).get(transaction['merchant_id'], 0)
        
        # Caractéristiques de communauté
        features['customer_community'] = graph_data.get('customer_community', {}).get(transaction['customer_id'], -1)
        features['merchant_community'] = graph_data.get('merchant_community', {}).get(transaction['merchant_id'], -1)
        
        # Caractéristiques de cycle
        features['is_in_cycle'] = 1 if self.detect_cycles(transaction, graph_data) else 0
        
        return features
    
    def detect_cycles(self, transaction: dict, graph_data: dict) -> bool:
        """Détecte si la transaction fait partie d'un cycle suspect"""
        # Implémentation simplifiée de détection de cycle
        customer_id = transaction['customer_id']
        merchant_id = transaction['merchant_id']
        
        # Vérifier si le merchant a transigé avec des clients qui ont aussi transigé avec ce customer
        customer_connections = graph_data.get('customer_connections', {}).get(customer_id, set())
        merchant_connections = graph_data.get('merchant_connections', {}).get(merchant_id, set())
        
        return len(customer_connections.intersection(merchant_connections)) > 2
\end{minted}

\section{Algorithmes de Machine Learning pour la Détection de Fraude}

\subsection{Random Forest}

\begin{minted}{python}
from sklearn.ensemble import RandomForestClassifier
from sklearn.model_selection import cross_val_score
import numpy as np

class RandomForestFraudDetector:
    def __init__(self, n_estimators: int = 100, max_depth: int = 10):
        self.model = RandomForestClassifier(
            n_estimators=n_estimators,
            max_depth=max_depth,
            random_state=42,
            class_weight='balanced',
            n_jobs=-1,
            min_samples_split=10,
            min_samples_leaf=5
        )
        self.feature_importance = None
    
    def train(self, X_train, y_train):
        """Entraîne le modèle Random Forest"""
        self.model.fit(X_train, y_train)
        self.feature_importance = dict(zip(
            X_train.columns,
            self.model.feature_importances_
        ))
    
    def predict(self, X):
        """Prédit les classes"""
        return self.model.predict(X)
    
    def predict_proba(self, X):
        """Prédit les probabilités"""
        return self.model.predict_proba(X)
    
    def get_feature_importance(self, top_n: int = 10):
        """Retourne les caractéristiques les plus importantes"""
        if self.feature_importance is None:
            return None
        
        sorted_features = sorted(
            self.feature_importance.items(),
            key=lambda x: x[1],
            reverse=True
        )
        
        return sorted_features[:top_n]
    
    def evaluate(self, X_test, y_test):
        """Évalue le modèle avec plusieurs métriques"""
        from sklearn.metrics import (
            accuracy_score, precision_score, recall_score,
            f1_score, roc_auc_score, confusion_matrix
        )
        
        y_pred = self.predict(X_test)
        y_proba = self.predict_proba(X_test)[:, 1]
        
        metrics = {
            'accuracy': accuracy_score(y_test, y_pred),
            'precision': precision_score(y_test, y_pred),
            'recall': recall_score(y_test, y_pred),
            'f1_score': f1_score(y_test, y_pred),
            'roc_auc': roc_auc_score(y_test, y_proba),
            'confusion_matrix': confusion_matrix(y_test, y_pred)
        }
        
        return metrics
\end{minted}

\subsection{Gradient Boosting (XGBoost/LightGBM)}

\begin{minted}{python}
import xgboost as xgb
from sklearn.model_selection import train_test_split

class XGBoostFraudDetector:
    def __init__(self, scale_pos_weight: float = None):
        self.scale_pos_weight = scale_pos_weight
        self.model = None
        self.feature_importance = None
    
    def train(self, X_train, y_train, X_val=None, y_val=None):
        """Entraîne le modèle XGBoost"""
        # Calcul automatique du scale_pos_weight si non fourni
        if self.scale_pos_weight is None:
            neg_count = (y_train == 0).sum()
            pos_count = (y_train == 1).sum()
            self.scale_pos_weight = neg_count / pos_count
        
        params = {
            'objective': 'binary:logistic',
            'eval_metric': 'auc',
            'scale_pos_weight': self.scale_pos_weight,
            'max_depth': 6,
            'learning_rate': 0.1,
            'subsample': 0.8,
            'colsample_bytree': 0.8,
            'random_state': 42,
            'n_jobs': -1
        }
        
        dtrain = xgb.DMatrix(X_train, label=y_train)
        
        if X_val is not None and y_val is not None:
            dval = xgb.DMatrix(X_val, label=y_val)
            self.model = xgb.train(
                params,
                dtrain,
                num_boost_round=100,
                evals=[(dtrain, 'train'), (dval, 'val')],
                early_stopping_rounds=10,
                verbose_eval=10
            )
        else:
            self.model = xgb.train(params, dtrain, num_boost_round=100)
        
        self.feature_importance = self.model.get_score(importance_type='gain')
    
    def predict(self, X):
        """Prédit les classes"""
        dtest = xgb.DMatrix(X)
        proba = self.model.predict(dtest)
        return (proba > 0.5).astype(int)
    
    def predict_proba(self, X):
        """Prédit les probabilités"""
        dtest = xgb.DMatrix(X)
        return self.model.predict(dtest)
    
    def get_feature_importance(self, top_n: int = 10):
        """Retourne les caractéristiques les plus importantes"""
        if self.feature_importance is None:
            return None
        
        sorted_features = sorted(
            self.feature_importance.items(),
            key=lambda x: x[1],
            reverse=True
        )
        
        return sorted_features[:top_n]
    
    def save_model(self, filepath: str):
        """Sauvegarde le modèle"""
        self.model.save_model(filepath)
    
    def load_model(self, filepath: str):
        """Charge un modèle sauvegardé"""
        self.model = xgb.Booster()
        self.model.load_model(filepath)
        self.feature_importance = self.model.get_score(importance_type='gain')
\end{minted}

\subsection{AutoML pour la Détection de Fraude}

\begin{minted}{python}
from sklearn.ensemble import RandomForestClassifier, GradientBoostingClassifier
from sklearn.linear_model import LogisticRegression
from sklearn.svm import SVC
from sklearn.model_selection import cross_val_score
import numpy as np

class AutoMLFraudDetector:
    def __init__(self):
        self.models = {
            'random_forest': RandomForestClassifier(
                n_estimators=100, class_weight='balanced', random_state=42
            ),
            'gradient_boosting': GradientBoostingClassifier(
                n_estimators=100, random_state=42
            ),
            'logistic_regression': LogisticRegression(
                class_weight='balanced', random_state=42, max_iter=1000
            ),
            'svm': SVC(
                class_weight='balanced', probability=True, random_state=42
            )
        }
        self.best_model = None
        self.best_model_name = None
        self.best_score = 0
    
    def train_and_select(self, X_train, y_train, cv: int = 5):
        """Entraîne tous les modèles et sélectionne le meilleur"""
        results = {}
        
        for name, model in self.models.items():
            print(f"Training {name}...")
            
            try:
                scores = cross_val_score(
                    model, X_train, y_train,
                    cv=cv,
                    scoring='roc_auc',
                    n_jobs=-1
                )
                
                mean_score = scores.mean()
                std_score = scores.std()
                
                results[name] = {
                    'mean_score': mean_score,
                    'std_score': std_score,
                    'scores': scores
                }
                
                print(f"{name}: {mean_score:.4f} (+/- {std_score:.4f})")
                
                if mean_score > self.best_score:
                    self.best_score = mean_score
                    self.best_model_name = name
                    self.best_model = model
                    
            except Exception as e:
                print(f"Error training {name}: {e}")
                results[name] = {'error': str(e)}
        
        # Entraîner le meilleur modèle sur l'ensemble du dataset
        if self.best_model is not None:
            print(f"\nBest model: {self.best_model_name} with score {self.best_score:.4f}")
            self.best_model.fit(X_train, y_train)
        
        return results
    
    def predict(self, X):
        """Prédit avec le meilleur modèle"""
        if self.best_model is None:
            raise ValueError("No model trained. Call train_and_select first.")
        
        return self.best_model.predict(X)
    
    def predict_proba(self, X):
        """Prédit les probabilités avec le meilleur modèle"""
        if self.best_model is None:
            raise ValueError("No model trained. Call train_and_select first.")
        
        return self.best_model.predict_proba(X)
\end{minted}

\section{Explicabilité du Modèle (Model Explainability)}

\subsection{SHAP (SHapley Additive exPlanations)}

\begin{minted}{python}
import shap
import matplotlib.pyplot as plt

class ModelExplainer:
    def __init__(self, model, feature_names):
        self.model = model
        self.feature_names = feature_names
        self.explainer = None
    
    def setup_explainer(self, X_background):
        """Configure l'explainer SHAP"""
        try:
            # Pour les modèles tree-based (Random Forest, XGBoost)
            if hasattr(self.model, 'estimators_') or hasattr(self.model, 'feature_importances_'):
                self.explainer = shap.TreeExplainer(self.model)
            else:
                # Pour les autres modèles
                self.explainer = shap.KernelExplainer(
                    self.model.predict_proba,
                    X_background
                )
        except Exception as e:
            print(f"Error setting up explainer: {e}")
            self.explainer = shap.KernelExplainer(
                self.model.predict_proba,
                X_background
            )
    
    def explain_prediction(self, instance):
        """Explique une prédiction individuelle"""
        if self.explainer is None:
            raise ValueError("Explainer not set up. Call setup_explainer first.")
        
        shap_values = self.explainer.shap_values(instance)
        
        # Pour classification binaire, shap_values retourne [shap_values_class_0, shap_values_class_1]
        if isinstance(shap_values, list):
            shap_values = shap_values[1]  # On s'intéresse à la classe positive (fraude)
        
        explanation = {
            'prediction': self.model.predict([instance])[0],
            'probability': self.model.predict_proba([instance])[0][1],
            'shap_values': dict(zip(self.feature_names, shap_values)),
            'base_value': self.explainer.expected_value[1] if isinstance(self.explainer.expected_value, list) else self.explainer.expected_value
        }
        
        return explanation
    
    def get_global_feature_importance(self, X):
        """Retourne l'importance globale des caractéristiques"""
        if self.explainer is None:
            raise ValueError("Explainer not set up. Call setup_explainer first.")
        
        shap_values = self.explainer.shap_values(X)
        
        if isinstance(shap_values, list):
            shap_values = shap_values[1]
        
        # Calcul de l'importance moyenne absolue
        mean_abs_shap = np.abs(shap_values).mean(axis=0)
        
        importance = dict(zip(self.feature_names, mean_abs_shap))
        
        return sorted(importance.items(), key=lambda x: x[1], reverse=True)
    
    def plot_summary(self, X):
        """Génère le plot SHAP summary"""
        if self.explainer is None:
            raise ValueError("Explainer not set up. Call setup_explainer first.")
        
        shap_values = self.explainer.shap_values(X)
        
        if isinstance(shap_values, list):
            shap_values = shap_values[1]
        
        shap.summary_plot(shap_values, X, feature_names=self.feature_names)
    
    def plot_waterfall(self, instance):
        """Génère le plot waterfall pour une instance"""
        if self.explainer is None:
            raise ValueError("Explainer not set up. Call setup_explainer first.")
        
        shap_values = self.explainer.shap_values(instance)
        
        if isinstance(shap_values, list):
            shap_values = shap_values[1]
        
        shap.waterfall_plot(shap.Explanation(
            values=shap_values,
            base_values=self.explainer.expected_value[1] if isinstance(self.explainer.expected_value, list) else self.explainer.expected_value,
            data=instance,
            feature_names=self.feature_names
        ))
\end{minted}

\subsection{LIME (Local Interpretable Model-agnostic Explanations)}

\begin{minted}{python}
from lime.lime_tabular import LimeTabularExplainer

class LimeExplainer:
    def __init__(self, model, X_train, feature_names, class_names):
        self.model = model
        self.feature_names = feature_names
        self.class_names = class_names
        self.explainer = LimeTabularExplainer(
            X_train,
            feature_names=feature_names,
            class_names=class_names,
            mode='classification'
        )
    
    def explain_instance(self, instance, num_features: int = 10):
        """Explique une instance spécifique avec LIME"""
        explanation = self.explainer.explain_instance(
            instance,
            self.model.predict_proba,
            num_features=num_features
        )
        
        return explanation
    
    def get_explanation_dict(self, instance, num_features: int = 10):
        """Retourne l'explication sous forme de dictionnaire"""
        explanation = self.explain_instance(instance, num_features)
        
        exp_dict = {
            'prediction': self.model.predict([instance])[0],
            'probability': self.model.predict_proba([instance])[0][1],
            'local_features': []
        }
        
        for feature, weight in explanation.as_list():
            exp_dict['local_features'].append({
                'feature': feature,
                'weight': weight,
                'impact': 'positive' if weight > 0 else 'negative'
            })
        
        return exp_dict
    
    def plot_explanation(self, instance, num_features: int = 10):
        """Génère le plot LIME"""
        explanation = self.explain_instance(instance, num_features)
        explanation.as_pyplot_figure()
\end{minted}

\begin{bestpractice}
L'explicabilité des modèles est cruciale dans le domaine financier pour la conformité réglementaire et pour maintenir la confiance des utilisateurs. SHAP et LIME permettent de comprendre les décisions du modèle et d'identifier les biais potentiels.
\end{bestpractice}

% ==============================================================================
% CHAPITRE 7 : ÉTUDES DE CAS ET SCÉNARIOS D'USAGE
% ==============================================================================
\chapter{Études de Cas et Scénarios d'Usage}

\section{Scénario 1 : Intégration Multi-Banques pour une PME}

\subsection{Contexte et Besoins}

Une PME de taille moyenne dispose de comptes bancaires auprès de trois établissements différents pour diversifier ses risques et optimiser ses coûts bancaires. L'entreprise fait face aux défis suivants :

\begin{itemize}
    \item \textbf{Fragmentation des données} : Les informations bancaires sont dispersées sur plusieurs portails bancaires.
    \item \textbf{Processus manuels} : Les rapprochements bancaires nécessitent des exports manuels depuis chaque banque.
    \item \textbf{Délais de traitement} : Les décisions financières sont retardées par le temps de consolidation des données.
    \item \textbf{Risque d'erreurs} : Les manipulations manuelles augmentent le risque d'erreurs de saisie.
\end{itemize}

\subsection{Solution Technique}

\subsubsection{Architecture de Connexion}

L'implémentation utilise le Module F2 avec les composants suivants :

\begin{minted}{typescript}
interface PMEBankingIntegration {
  companyId: string;
  bankConnections: BankConnection[];
  syncSchedule: SyncSchedule;
  notificationPreferences: NotificationPreferences;
}

class PMEBankingOrchestrator {
  constructor(
    private bankHub: MultiBankingHub,
    private fraudDetector: FraudDetectionService,
    private notificationService: NotificationService
  ) {}

  async setupPMEIntegration(config: PMEBankingIntegration): Promise<IntegrationResult> {
    const results = [];
    
    // Création des connexions bancaires
    for (const bankConfig of config.bankConnections) {
      try {
        const connection = await this.bankHub.createConnection({
          companyId: config.companyId,
          bankId: bankConfig.bankId,
          connectionType: bankConfig.connectionType,
          credentials: bankConfig.credentials,
          consentValidityDays: 365
        });
        
        results.push({
          bankId: bankConfig.bankId,
          status: 'SUCCESS',
          connectionId: connection.connectionId
        });
        
      } catch (error) {
        results.push({
          bankId: bankConfig.bankId,
          status: 'FAILED',
          error: error.message
        });
      }
    }
    
    // Configuration de la synchronisation automatique
    await this.setupAutoSync(config.companyId, config.syncSchedule);
    
    // Configuration des notifications
    await this.notificationService.configureAlerts(
      config.companyId,
      config.notificationPreferences
    );
    
    return {
      companyId: config.companyId,
      connections: results,
      overallStatus: results.every(r => r.status === 'SUCCESS') ? 'SUCCESS' : 'PARTIAL'
    };
  }

  private async setupAutoSync(companyId: string, schedule: SyncSchedule): Promise<void> {
    // Configuration de la synchronisation quotidienne
    const cronJob = new CronJob(schedule.frequency, async () => {
      await this.syncAllBankAccounts(companyId);
    });
    
    cronJob.start();
  }

  async syncAllBankAccounts(companyId: string): Promise<SyncResult> {
    const connections = await this.bankHub.getActiveConnections(companyId);
    const syncResults = [];
    
    for (const connection of connections) {
      try {
        const syncResult = await this.bankHub.syncConnection(connection.id, {
          fromDate: new Date(Date.now() - 24 * 60 * 60 * 1000), // 24h
          includePending: true
        });
        
        // Analyse de fraude sur les nouvelles transactions
        for (const transaction of syncResult.transactions) {
          const fraudAnalysis = await this.fraudDetector.analyzeTransaction(transaction);
          
          if (fraudAnalysis.riskScore > 0.7) {
            await this.notificationService.sendHighRiskAlert({
              companyId,
              transactionId: transaction.id,
              riskScore: fraudAnalysis.riskScore,
              reasons: fraudAnalysis.reasons
            });
          }
        }
        
        syncResults.push({
          connectionId: connection.id,
          status: 'SUCCESS',
          transactionCount: syncResult.transactions.length
        });
        
      } catch (error) {
        syncResults.push({
          connectionId: connection.id,
          status: 'FAILED',
          error: error.message
        });
      }
    }
    
    return {
      companyId,
      syncResults,
      timestamp: new Date()
    };
  }
}
\end{minted}

\subsubsection{Tableau de Bord de Consolidation}

\begin{minted}{typescript}
interface ConsolidatedDashboard {
  companyId: string;
  overview: FinancialOverview;
  bankBreakdown: BankBreakdown[];
  cashFlow: CashFlowAnalysis;
  alerts: Alert[];
}

class PMEDashboardService {
  constructor(
    private bankHub: MultiBankingHub,
    private analyticsService: AnalyticsService
  ) {}

  async generateDashboard(companyId: string): Promise<ConsolidatedDashboard> {
    // Récupération des données de toutes les banques
    const allTransactions = await this.bankHub.getAllTransactions(companyId, {
      startDate: new Date(Date.now() - 30 * 24 * 60 * 60 * 1000), // 30 jours
      endDate: new Date()
    });

    // Calcul des vues globales
    const overview = this.calculateOverview(allTransactions);
    
    // Analyse par banque
    const bankBreakdown = await this.analyzeByBank(companyId, allTransactions);
    
    // Analyse de trésorerie
    const cashFlow = this.analyzeCashFlow(allTransactions);
    
    // Récupération des alertes
    const alerts = await this.getRecentAlerts(companyId);

    return {
      companyId,
      overview,
      bankBreakdown,
      cashFlow,
      alerts
    };
  }

  private calculateOverview(transactions: Transaction[]): FinancialOverview {
    const totalBalance = transactions.reduce((sum, t) => sum + t.amount, 0);
    const income = transactions.filter(t => t.amount > 0).reduce((sum, t) => sum + t.amount, 0);
    const expenses = transactions.filter(t => t.amount < 0).reduce((sum, t) => sum + Math.abs(t.amount), 0);
    
    return {
      totalBalance,
      totalIncome: income,
      totalExpenses: expenses,
      netCashFlow: income - expenses,
      transactionCount: transactions.length,
      averageTransactionAmount: totalBalance / transactions.length
    };
  }

  private async analyzeByBank(companyId: string, transactions: Transaction[]): Promise<BankBreakdown[]> {
    const connections = await this.bankHub.getActiveConnections(companyId);
    const breakdown = [];
    
    for (const connection of connections) {
      const bankTransactions = transactions.filter(t => t.connectionId === connection.id);
      
      breakdown.push({
        bankId: connection.bankId,
        bankName: connection.bankName,
        balance: bankTransactions.reduce((sum, t) => sum + t.amount, 0),
        transactionCount: bankTransactions.length,
        lastSync: connection.lastSyncDate
      });
    }
    
    return breakdown.sort((a, b) => b.balance - a.balance);
  }

  private analyzeCashFlow(transactions: Transaction[]): CashFlowAnalysis {
    // Analyse des tendances de trésorerie
    const dailyBalances = this.calculateDailyBalances(transactions);
    const trend = this.calculateTrend(dailyBalances);
    
    return {
      currentBalance: dailyBalances[dailyBalances.length - 1],
      averageBalance: dailyBalances.reduce((a, b) => a + b) / dailyBalances.length,
      trend: trend > 0 ? 'INCREASING' : trend < 0 ? 'DECREASING' : 'STABLE',
      trendPercentage: Math.abs(trend) * 100,
      projectedBalance: this.projectBalance(dailyBalances, 7) // Projection 7 jours
    };
  }
}
\end{minted}

\subsection{Résultats et Bénéfices}

L'implémentation a permis à la PME de réaliser les améliorations suivantes :

\begin{itemize}
    \item \textbf{Gain de temps} : Réduction de 80\% du temps consacré aux rapprochements bancaires.
    \item \textbf{Meilleure visibilité} : Vue consolidée en temps réel de la position de trésorerie.
    \item \textbf{Réduction des erreurs} : Élimination des erreurs de saisie manuelle.
    \item \textbf{Détection proactive} : Identification automatique des transactions suspectes.
    \item \textbf{Prise de décision} : Accélération des décisions financières grâce aux données en temps réel.
\end{itemize}

\section{Scénario 2 : Détection de Fraude Complexes par Graph Analysis}

\subsection{Contexte du Cas}

Un réseau de blanchiment d'argent sophistiqué utilise des transactions en cascade pour dissimuler l'origine des fonds. Les patterns traditionnels basés sur des règles ne parviennent pas à détecter ce type de fraude complexe.

\subsection{Analyse de Graphe}

\subsubsection{Construction du Graphe de Transactions}

\begin{minted}{python}
import networkx as nx
from typing import List, Dict

class FraudGraphAnalyzer:
    def __init__(self):
        self.graph = nx.DiGraph()
        self.node_attributes = {}
        self.edge_attributes = {}
    
    def build_transaction_graph(self, transactions: List[Dict]) -> nx.DiGraph:
        """Construit un graphe orienté des transactions"""
        self.graph = nx.DiGraph()
        
        for tx in transactions:
            sender = tx['sender_account']
            receiver = tx['receiver_account']
            amount = tx['amount']
            timestamp = tx['timestamp']
            
            # Ajout des nœuds s'ils n'existent pas
            if not self.graph.has_node(sender):
                self.graph.add_node(sender, 
                                   account_type=tx.get('sender_type', 'unknown'),
                                   risk_score=tx.get('sender_risk', 0))
            
            if not self.graph.has_node(receiver):
                self.graph.add_node(receiver,
                                   account_type=tx.get('receiver_type', 'unknown'),
                                   risk_score=tx.get('receiver_risk', 0))
            
            # Ajout de l'arête transactionnelle
            edge_key = (sender, receiver)
            if self.graph.has_edge(*edge_key):
                # Agrégation si l'arête existe déjà
                self.graph[edge_key[0]][edge_key[1]]['total_amount'] += amount
                self.graph[edge_key[0]][edge_key[1]]['transaction_count'] += 1
                self.graph[edge_key[0]][edge_key[1]]['transactions'].append({
                    'amount': amount,
                    'timestamp': timestamp,
                    'transaction_id': tx['id']
                })
            else:
                self.graph.add_edge(sender, receiver,
                                  total_amount=amount,
                                  transaction_count=1,
                                  transactions=[{
                                      'amount': amount,
                                      'timestamp': timestamp,
                                      'transaction_id': tx['id']
                                  }])
        
        return self.graph
    
    def detect_circular_patterns(self, min_cycle_length: int = 3, max_cycle_length: int = 10) -> List[Dict]:
        """Détecte les patterns circulaires suspectes"""
        suspicious_cycles = []
        
        # Recherche des cycles simples
        cycles = list(nx.simple_cycles(self.graph))
        
        for cycle in cycles:
            if min_cycle_length <= len(cycle) <= max_cycle_length:
                cycle_analysis = self.analyze_cycle(cycle)
                
                if cycle_analysis['suspicion_score'] > 0.7:
                    suspicious_cycles.append(cycle_analysis)
        
        return suspicious_cycles
    
    def analyze_cycle(self, cycle: List[str]) -> Dict:
        """Analyse un cycle de transactions"""
        total_amount = 0
        transaction_count = 0
        time_span = 0
        amounts = []
        
        for i in range(len(cycle)):
            current = cycle[i]
            next_node = cycle[(i + 1) % len(cycle)]
            
            if self.graph.has_edge(current, next_node):
                edge_data = self.graph[current][next_node]
                total_amount += edge_data['total_amount']
                transaction_count += edge_data['transaction_count']
                amounts.extend([t['amount'] for t in edge_data['transactions']])
        
        # Calcul du score de suspicion
        suspicion_score = self.calculate_cycle_suspicion(
            cycle, total_amount, transaction_count, amounts
        )
        
        return {
            'cycle': cycle,
            'length': len(cycle),
            'total_amount': total_amount,
            'transaction_count': transaction_count,
            'average_amount': total_amount / transaction_count if transaction_count > 0 else 0,
            'amount_variance': np.var(amounts) if amounts else 0,
            'suspicion_score': suspicion_score,
            'risk_factors': self.identify_risk_factors(cycle, total_amount, amounts)
        }
    
    def calculate_cycle_suspicion(self, cycle: List[str], total_amount: float, 
                                   transaction_count: int, amounts: List[float]) -> float:
        """Calcule un score de suspicion pour un cycle"""
        score = 0.0
        
        # Facteur 1: Nombre de transactions élevé
        if transaction_count > 10:
            score += 0.3
        
        # Facteur 2: Montants similaires (pattern de structuration)
        if len(amounts) > 1:
            cv = np.std(amounts) / np.mean(amounts) if np.mean(amounts) > 0 else 0
            if cv < 0.2:  # Coefficient de variation faible
                score += 0.3
        
        # Facteur 3: Montant total élevé
        if total_amount > 10000:
            score += 0.2
        
        # Facteur 4: Longueur du cycle
        if len(cycle) > 5:
            score += 0.2
        
        return min(score, 1.0)
    
    def identify_risk_factors(self, cycle: List[str], total_amount: float, 
                             amounts: List[float]) -> List[str]:
        """Identifie les facteurs de risque spécifiques"""
        risk_factors = []
        
        if len(amounts) > 1:
            cv = np.std(amounts) / np.mean(amounts) if np.mean(amounts) > 0 else 0
            if cv < 0.1:
                risk_factors.append("STRUCTURED_AMOUNTS")
        
        if total_amount > 50000:
            risk_factors.append("HIGH_VOLUME")
        
        if len(cycle) > 7:
            risk_factors.append("COMPLEX_STRUCTURE")
        
        # Vérification des comptes à risque
        for node in cycle:
            if self.graph.nodes[node].get('risk_score', 0) > 0.7:
                risk_factors.append("HIGH_RISK_ACCOUNT")
                break
        
        return risk_factors
    
    def detect_money_laundering_patterns(self) -> List[Dict]:
        """Détecte des patterns de blanchiment d'argent"""
        patterns = []
        
        # Pattern 1: Fan-out (un compte vers plusieurs)
        fan_out_accounts = self.detect_fan_out(threshold=5)
        patterns.extend(fan_out_accounts)
        
        # Pattern 2: Fan-in (plusieurs comptes vers un)
        fan_in_accounts = self.detect_fan_in(threshold=5)
        patterns.extend(fan_in_accounts)
        
        # Pattern 3: Short cycles
        short_cycles = self.detect_circular_patterns(min_cycle_length=2, max_cycle_length=4)
        patterns.extend(short_cycles)
        
        return patterns
    
    def detect_fan_out(self, threshold: int = 5) -> List[Dict]:
        """Détecte les comptes qui envoient à de nombreux destinataires"""
        fan_out_patterns = []
        
        for node in self.graph.nodes():
            out_degree = self.graph.out_degree(node)
            if out_degree >= threshold:
                total_amount = sum(
                    self.graph[node][neighbor]['total_amount']
                    for neighbor in self.graph.successors(node)
                )
                
                fan_out_patterns.append({
                    'pattern_type': 'FAN_OUT',
                    'account': node,
                    'out_degree': out_degree,
                    'total_amount': total_amount,
                    'destinations': list(self.graph.successors(node))
                })
        
        return fan_out_patterns
    
    def detect_fan_in(self, threshold: int = 5) -> List[Dict]:
        """Détecte les comptes qui reçoivent de nombreux expéditeurs"""
        fan_in_patterns = []
        
        for node in self.graph.nodes():
            in_degree = self.graph.in_degree(node)
            if in_degree >= threshold:
                total_amount = sum(
                    self.graph[neighbor][node]['total_amount']
                    for neighbor in self.graph.predecessors(node)
                )
                
                fan_in_patterns.append({
                    'pattern_type': 'FAN_IN',
                    'account': node,
                    'in_degree': in_degree,
                    'total_amount': total_amount,
                    'sources': list(self.graph.predecessors(node))
                })
        
        return fan_in_patterns
\end{minted}

\subsection{Intégration avec le Système de Détection de Fraude}

\begin{minted}{python}
class HybridFraudDetectionPipeline:
    def __init__(self):
        self.rule_engine = RuleBasedEngine()
        self.ml_detector = MachineLearningDetector()
        self.graph_analyzer = FraudGraphAnalyzer()
        self.alert_system = AlertSystem()
    
    def analyze_transaction(self, transaction: Dict) -> Dict:
        """Analyse complète d'une transaction avec approche hybride"""
        results = {
            'transaction_id': transaction['id'],
            'timestamp': transaction['timestamp'],
            'rule_based_analysis': None,
            'ml_analysis': None,
            'graph_analysis': None,
            'overall_risk_score': 0.0,
            'risk_level': 'LOW',
            'alert_triggered': False
        }
        
        # 1. Analyse basée sur les règles
        rule_result = self.rule_engine.analyze(transaction)
        results['rule_based_analysis'] = rule_result
        
        # 2. Analyse Machine Learning
        ml_result = self.ml_detector.predict(transaction)
        results['ml_analysis'] = ml_result
        
        # 3. Analyse de graphe (si contexte disponible)
        if hasattr(transaction, 'related_transactions'):
            self.graph_analyzer.build_transaction_graph(
                transaction['related_transactions']
            )
            graph_result = self.graph_analyzer.detect_money_laundering_patterns()
            results['graph_analysis'] = graph_result
        
        # 4. Calcul du score de risque combiné
        results['overall_risk_score'] = self.calculate_combined_risk(
            rule_result, ml_result, results.get('graph_analysis')
        )
        
        # 5. Détermination du niveau de risque
        results['risk_level'] = self.determine_risk_level(results['overall_risk_score'])
        
        # 6. Déclenchement d'alerte si nécessaire
        if results['overall_risk_score'] > 0.7:
            results['alert_triggered'] = True
            self.alert_system.send_alert(results)
        
        return results
    
    def calculate_combined_risk(self, rule_result: Dict, ml_result: Dict, 
                                graph_result: List[Dict] = None) -> float:
        """Calcule un score de risque combiné"""
        # Pondération des différentes analyses
        rule_weight = 0.3
        ml_weight = 0.5
        graph_weight = 0.2
        
        rule_score = rule_result.get('risk_score', 0)
        ml_score = ml_result.get('fraud_probability', 0)
        
        graph_score = 0
        if graph_result and len(graph_result) > 0:
            # Score basé sur le nombre de patterns suspects
            graph_score = min(len(graph_result) * 0.2, 1.0)
        
        combined_score = (
            rule_weight * rule_score +
            ml_weight * ml_score +
            graph_weight * graph_score
        )
        
        return combined_score
    
    def determine_risk_level(self, risk_score: float) -> str:
        """Détermine le niveau de risque basé sur le score"""
        if risk_score >= 0.8:
            return 'CRITICAL'
        elif risk_score >= 0.6:
            return 'HIGH'
        elif risk_score >= 0.4:
            return 'MEDIUM'
        else:
            return 'LOW'
\end{minted}

\begin{infobox}
L'approche hybride combinant règles, Machine Learning et analyse de graphe permet de détecter des patterns de fraude complexes qui échapperaient aux méthodes traditionnelles. L'analyse de graphe est particulièrement efficace pour détecter les réseaux de blanchiment d'argent.
\end{infobox}

\section{Scénario 3 : Rapprochement Bancaire Automatisé}

\subsection{Processus de Rapprochement Traditionnel}

Le rapprochement bancaire traditionnel implique les étapes manuelles suivantes :

\begin{enumerate}
    \item Export des relevés bancaires depuis chaque interface bancaire
    \item Export des écritures comptables depuis le système comptable
    \item Rapprochement manuel ligne par ligne
    \item Identification des écarts
    \item Écritures de régularisation
    \item Validation et archivage
\end{enumerate}

\subsection{Solution Automatisée}

\begin{minted}{typescript}
interface ReconciliationConfig {
  companyId: string;
  bankConnectionIds: string[];
  accountingSystemId: string;
  reconciliationRules: ReconciliationRule[];
  tolerance: ToleranceConfig;
}

interface ReconciliationResult {
  companyId: string;
  period: DateRange;
  totalBankTransactions: number;
  totalAccountingEntries: number;
  matchedItems: MatchedItem[];
  unmatchedBankItems: UnmatchedItem[];
  unmatchedAccountingItems: UnmatchedItem[];
  autoReconciled: number;
  manualReviewRequired: number;
  reconciliationRate: number;
}

class AutomatedReconciliationService {
  constructor(
    private bankHub: MultiBankingHub,
    private accountingIntegration: AccountingIntegration,
    private matchingEngine: MatchingEngine
  ) {}

  async performReconciliation(config: ReconciliationConfig): Promise<ReconciliationResult> {
    // 1. Récupération des données bancaires
    const bankData = await this.fetchBankData(config);
    
    // 2. Récupération des données comptables
    const accountingData = await this.fetchAccountingData(config);
    
    // 3. Normalisation des données
    const normalizedBank = this.normalizeBankData(bankData);
    const normalizedAccounting = this.normalizeAccountingData(accountingData);
    
    // 4. Matching intelligent
    const matchingResults = await this.matchingEngine.match(
      normalizedBank,
      normalizedAccounting,
      config.reconciliationRules,
      config.tolerance
    );
    
    // 5. Classification des résultats
    const classifiedResults = this.classifyResults(matchingResults);
    
    // 6. Génération des écritures de régularisation automatiques
    const autoAdjustments = await this.generateAutoAdjustments(
      classifiedResults.autoMatchable
    );
    
    // 7. Calcul des statistiques
    const statistics = this.calculateStatistics(
      normalizedBank,
      normalizedAccounting,
      classifiedResults
    );
    
    return {
      companyId: config.companyId,
      period: config.period,
      totalBankTransactions: normalizedBank.length,
      totalAccountingEntries: normalizedAccounting.length,
      matchedItems: classifiedResults.matched,
      unmatchedBankItems: classifiedResults.unmatchedBank,
      unmatchedAccountingItems: classifiedResults.unmatchedAccounting,
      autoReconciled: autoAdjustments.length,
      manualReviewRequired: classifiedResults.requireManualReview.length,
      reconciliationRate: statistics.reconciliationRate
    };
  }

  private async fetchBankData(config: ReconciliationConfig): Promise<BankTransaction[]> {
    const allTransactions = [];
    
    for (const connectionId of config.bankConnectionIds) {
      const transactions = await this.bankHub.getTransactions(connectionId, {
        startDate: config.period.startDate,
        endDate: config.period.endDate
      });
      allTransactions.push(...transactions);
    }
    
    return allTransactions;
  }

  private async fetchAccountingData(config: ReconciliationConfig): Promise<AccountingEntry[]> {
    return await this.accountingIntegration.getEntries(config.accountingSystemId, {
      startDate: config.period.startDate,
      endDate: config.period.endDate,
      accountTypes: ['BANK', 'CASH']
    });
  }

  private normalizeBankData(transactions: BankTransaction[]): NormalizedTransaction[] {
    return transactions.map(tx => ({
      id: tx.id,
      date: tx.transactionDate,
      amount: tx.amount,
      description: tx.description,
      reference: tx.reference,
      counterparty: tx.counterparty,
      category: this.categorizeTransaction(tx),
      normalizedAmount: Math.abs(tx.amount),
      normalizedDate: this.normalizeDate(tx.transactionDate)
    }));
  }

  private normalizeAccountingData(entries: AccountingEntry[]): NormalizedEntry[] {
    return entries.map(entry => ({
      id: entry.id,
      date: entry.entryDate,
      amount: entry.amount,
      description: entry.description,
      reference: entry.reference,
      account: entry.accountNumber,
      normalizedAmount: Math.abs(entry.amount),
      normalizedDate: this.normalizeDate(entry.entryDate)
    }));
  }

  private categorizeTransaction(transaction: BankTransaction): string {
    // Catégorisation intelligente basée sur la description et le montant
    const description = transaction.description.toLowerCase();
    
    if (description.includes('salari') || description.includes('paye')) {
      return 'PAYROLL';
    } else if (description.includes('fournisseur') || description.includes('facture')) {
      return 'SUPPLIER';
    } else if (description.includes('client') || description.includes('vente')) {
      return 'CUSTOMER';
    } else if (description.includes('taxe') || description.includes('impôt')) {
      return 'TAX';
    } else {
      return 'OTHER';
    }
  }

  private normalizeDate(date: Date): string {
    return date.toISOString().split('T')[0]; // YYYY-MM-DD
  }

  private classifyResults(matchingResults: MatchingResult): ClassifiedResults {
    return {
      matched: matchingResults.exactMatches,
      autoMatchable: matchingResults.fuzzyMatches.filter(m => m.confidence > 0.9),
      requireManualReview: matchingResults.fuzzyMatches.filter(m => m.confidence <= 0.9),
      unmatchedBank: matchingResults.unmatchedBank,
      unmatchedAccounting: matchingResults.unmatchedAccounting
    };
  }

  private calculateStatistics(
    bankData: NormalizedTransaction[],
    accountingData: NormalizedEntry[],
    classified: ClassifiedResults
  ): ReconciliationStatistics {
    const totalItems = bankData.length + accountingData.length;
    const matchedItems = classified.matched.length + classified.autoMatchable.length;
    
    return {
      reconciliationRate: matchedItems / totalItems,
      bankCoverage: (classified.matched.length + classified.unmatchedBank.length) / bankData.length,
      accountingCoverage: (classified.matched.length + classified.unmatchedAccounting.length) / accountingData.length,
      averageConfidence: this.calculateAverageConfidence(classified.autoMatchable)
    };
  }

  private calculateAverageConfidence(matches: FuzzyMatch[]): number {
    if (matches.length === 0) return 0;
    const totalConfidence = matches.reduce((sum, m) => sum + m.confidence, 0);
    return totalConfidence / matches.length;
  }
}

class MatchingEngine {
  async match(
    bankData: NormalizedTransaction[],
    accountingData: NormalizedEntry[],
    rules: ReconciliationRule[],
    tolerance: ToleranceConfig
  ): Promise<MatchingResult> {
    const exactMatches = [];
    const fuzzyMatches = [];
    const unmatchedBank = [...bankData];
    const unmatchedAccounting = [...accountingData];

    // 1. Matching exact par référence
    const referenceMatches = this.matchByReference(bankData, accountingData);
    exactMatches.push(...referenceMatches.matches);
    
    // Retirer les éléments matchés
    this.removeMatchedItems(unmatchedBank, unmatchedAccounting, referenceMatches.matchedIds);

    // 2. Matching flou par montant et date
    const fuzzyMatchResults = this.fuzzyMatchByAmountAndDate(
      unmatchedBank,
      unmatchedAccounting,
      tolerance
    );
    fuzzyMatches.push(...fuzzyMatchResults.matches);
    
    // Retirer les éléments matchés
    this.removeMatchedItems(unmatchedBank, unmatchedAccounting, fuzzyMatchResults.matchedIds);

    // 3. Matching par description et montant
    const descriptionMatches = this.matchByDescriptionAndAmount(
      unmatchedBank,
      unmatchedAccounting,
      tolerance
    );
    fuzzyMatches.push(...descriptionMatches.matches);

    return {
      exactMatches,
      fuzzyMatches,
      unmatchedBank,
      unmatchedAccounting
    };
  }

  private matchByReference(bankData: NormalizedTransaction[], accountingData: NormalizedEntry[]): MatchResult {
    const matches = [];
    const matchedIds = { bank: new Set(), accounting: new Set() };

    const accountingByReference = new Map(
      accountingData.map(entry => [entry.reference, entry])
    );

    for (const bankTx of bankData) {
      if (accountingByReference.has(bankTx.reference)) {
        const accountingEntry = accountingByReference.get(bankTx.reference);
        
        if (Math.abs(bankTx.normalizedAmount - accountingEntry.normalizedAmount) < 0.01) {
          matches.push({
            bankItem: bankTx,
            accountingItem: accountingEntry,
            matchType: 'EXACT_REFERENCE',
            confidence: 1.0
          });
          
          matchedIds.bank.add(bankTx.id);
          matchedIds.accounting.add(accountingEntry.id);
        }
      }
    }

    return { matches, matchedIds };
  }

  private fuzzyMatchByAmountAndDate(
    bankData: NormalizedTransaction[],
    accountingData: NormalizedEntry[],
    tolerance: ToleranceConfig
  ): MatchResult {
    const matches = [];
    const matchedIds = { bank: new Set(), accounting: new Set() };

    for (const bankTx of bankData) {
      for (const accountingEntry of accountingData) {
        const amountDiff = Math.abs(bankTx.normalizedAmount - accountingEntry.normalizedAmount);
        const dateDiff = Math.abs(
          new Date(bankTx.normalizedDate).getTime() - new Date(accountingEntry.normalizedDate).getTime()
        );
        const daysDiff = dateDiff / (1000 * 60 * 60 * 24);

        if (amountDiff <= tolerance.amountTolerance && daysDiff <= tolerance.dateToleranceDays) {
          const confidence = this.calculateFuzzyConfidence(amountDiff, daysDiff, tolerance);
          
          if (confidence > 0.7) {
            matches.push({
              bankItem: bankTx,
              accountingItem: accountingEntry,
              matchType: 'FUZZY_AMOUNT_DATE',
              confidence
            });
            
            matchedIds.bank.add(bankTx.id);
            matchedIds.accounting.add(accountingEntry.id);
          }
        }
      }
    }

    return { matches, matchedIds };
  }

  private calculateFuzzyConfidence(amountDiff: number, daysDiff: number, tolerance: ToleranceConfig): number {
    const amountScore = 1 - (amountDiff / tolerance.amountTolerance);
    const dateScore = 1 - (daysDiff / tolerance.dateToleranceDays);
    
    return (amountScore * 0.6) + (dateScore * 0.4);
  }

  private removeMatchedItems(
    bankData: NormalizedTransaction[],
    accountingData: NormalizedEntry[],
    matchedIds: { bank: Set<string>, accounting: Set<string> }
  ): void {
    // Filtrer les éléments déjà matchés
    const bankFilter = (item: NormalizedTransaction) => !matchedIds.bank.has(item.id);
    const accountingFilter = (item: NormalizedEntry) => !matchedIds.accounting.has(item.id);
    
    bankData.splice(0, bankData.length, ...bankData.filter(bankFilter));
    accountingData.splice(0, accountingData.length, ...accountingData.filter(accountingFilter));
  }
}
\end{minted}

\subsection{Bénéfices de l'Automatisation}

\begin{itemize}
    \item \textbf{Réduction du temps} : Le rapprochement automatique réduit le temps de traitement de 90\%.
    \item \textbf{Précision améliorée} : Réduction significative des erreurs humaines.
    \item \textbf{Détection précoce} : Identification rapide des écarts et anomalies.
    \item \textbf{Traçabilité} : Historique complet des opérations de rapprochement.
    \item \textbf{Scalabilité} : Capacité à traiter des volumes croissants de transactions.
\end{itemize}

\begin{bestpractice}
L'automatisation du rapprochement bancaire ne doit pas éliminer complètement l'intervention humaine. Les éléments non conciliés et les matchs avec faible confiance doivent toujours faire l'objet d'une revue manuelle pour garantir l'exactitude comptable.
\end{bestpractice}

% ==============================================================================
% CHAPITRE 8 : MONITORING ET OBSERVABILITÉ
% ==============================================================================
\chapter{Monitoring et Observabilité}

\section{Fondamentaux de l'Observabilité}

\subsection{Piliers de l'Observabilité}

L'observabilité moderne repose sur trois piliers fondamentaux :

\begin{itemize}
    \item \textbf{Metrics (Métriques)} : Données numériques agrégées sur le temps (ex: taux d'erreur, latence).
    \item \textbf{Logs (Journaux)} : Enregistrements d'événements discrets avec contexte (ex: messages d'erreur).
    \item \textbf{Traces (Traces)} : Suivi des requêtes à travers les services distribués (ex: distributed tracing).
\end{itemize}

\subsection{Objectifs du Monitoring Financier}

Le monitoring dans le contexte financier poursuit des objectifs spécifiques :

\begin{itemize}
    \item \textbf{Disponibilité} : Garantir que les services critiques sont toujours accessibles.
    \item \textbf{Performance} : Maintenir des temps de réponse acceptables pour les opérations financières.
    \item \textbf{Conformité} : Surveiller le respect des SLA et des exigences réglementaires.
    \item \textbf{Sécurité} : Détecter les activités suspectes et les tentatives d'intrusion.
    \item \textbf{Qualité des données} : Identifier les incohérences et les anomalies de données.
\end{itemize}

\section{Métriques et Alertes}

\subsection{Métriques de Base (RED Method)}

La méthode RED (Rate, Errors, Duration) fournit un cadre simple pour les métriques clés :

\begin{minted}{typescript}
import { Counter, Histogram, Gauge, Registry } from 'prom-client';

class REDMetrics {
  private requestCounter: Counter;
  private errorCounter: Counter;
  private durationHistogram: Histogram;
  private activeConnectionsGauge: Gauge;

  constructor(register: Registry) {
    this.requestCounter = new Counter({
      name: 'http_requests_total',
      help: 'Total number of HTTP requests',
      labelNames: ['method', 'route', 'status'],
      registers: [register]
    });

    this.errorCounter = new Counter({
      name: 'http_errors_total',
      help: 'Total number of HTTP errors',
      labelNames: ['method', 'route', 'error_type'],
      registers: [register]
    });

    this.durationHistogram = new Histogram({
      name: 'http_request_duration_seconds',
      help: 'Duration of HTTP requests in seconds',
      labelNames: ['method', 'route'],
      buckets: [0.005, 0.01, 0.025, 0.05, 0.1, 0.25, 0.5, 1, 2.5, 5, 10],
      registers: [register]
    });

    this.activeConnectionsGauge = new Gauge({
      name: 'active_bank_connections',
      help: 'Number of active bank connections',
      labelNames: ['bank_id', 'connection_type'],
      registers: [register]
    });
  }

  recordRequest(method: string, route: string, status: number, duration: number): void {
    this.requestCounter.inc({ method, route, status: status.toString() });
    
    if (status >= 400) {
      this.errorCounter.inc({ 
        method, 
        route, 
        error_type: status >= 500 ? 'server_error' : 'client_error' 
      });
    }
    
    this.durationHistogram.observe({ method, route }, duration / 1000);
  }

  updateActiveConnections(bankId: string, connectionType: string, count: number): void {
    this.activeConnectionsGauge.set({ bank_id: bankId, connection_type: connectionType }, count);
  }
}
\end{minted}

\subsection{Métriques Spécifiques aux Services Financiers}

\begin{minted}{typescript}
class FinancialMetrics {
  private bankSyncCounter: Counter;
  private fraudDetectionCounter: Counter;
  private transactionVolumeGauge: Gauge;
  private reconciliationRateGauge: Gauge;
  private apiLatencyHistogram: Histogram;

  constructor(register: Registry) {
    // Métriques pour le Multi Banking Hub
    this.bankSyncCounter = new Counter({
      name: 'bank_sync_operations_total',
      help: 'Total number of bank synchronization operations',
      labelNames: ['bank_id', 'status', 'connection_type'],
      registers: [register]
    });

    // Métriques pour la détection de fraude
    this.fraudDetectionCounter = new Counter({
      name: 'fraud_detections_total',
      help: 'Total number of fraud detections',
      labelNames: ['detection_type', 'severity', 'ml_confidence_range'],
      registers: [register]
    });

    // Métriques de volume de transactions
    this.transactionVolumeGauge = new Gauge({
      name: 'transaction_volume',
      help: 'Current transaction volume',
      labelNames: ['currency', 'transaction_type'],
      registers: [register]
    });

    // Métriques de taux de rapprochement
    this.reconciliationRateGauge = new Gauge({
      name: 'reconciliation_rate',
      help: 'Current reconciliation rate',
      labelNames: ['company_id', 'period'],
      registers: [register]
    });

    // Métriques de latence API externe
    this.apiLatencyHistogram = new Histogram({
      name: 'external_api_latency_seconds',
      help: 'Latency of external API calls',
      labelNames: ['api_name', 'operation'],
      buckets: [0.1, 0.5, 1, 2, 5, 10, 30, 60],
      registers: [register]
    });
  }

  recordBankSync(bankId: string, status: string, connectionType: string): void {
    this.bankSyncCounter.inc({ bank_id: bankId, status, connection_type });
  }

  recordFraudDetection(detectionType: string, severity: string, confidence: number): void {
    const confidenceRange = this.getConfidenceRange(confidence);
    this.fraudDetectionCounter.inc({ 
      detection_type: detectionType, 
      severity,
      ml_confidence_range: confidenceRange 
    });
  }

  updateTransactionVolume(currency: string, transactionType: string, volume: number): void {
    this.transactionVolumeGauge.set({ currency, transaction_type: transactionType }, volume);
  }

  updateReconciliationRate(companyId: string, period: string, rate: number): void {
    this.reconciliationRateGauge.set({ company_id: companyId, period }, rate);
  }

  recordExternalApiLatency(apiName: string, operation: string, latency: number): void {
    this.apiLatencyHistogram.observe({ api_name: apiName, operation }, latency / 1000);
  }

  private getConfidenceRange(confidence: number): string {
    if (confidence >= 0.9) return 'high';
    if (confidence >= 0.7) return 'medium';
    return 'low';
  }
}
\end{minted}

\subsection{Configuration des Alertes}

\begin{minted}{typescript}
interface AlertRule {
  name: string;
  condition: string;
  severity: 'critical' | 'warning' | 'info';
  duration: number; // en secondes
  annotations: Record<string, string>;
  labels: Record<string, string>;
}

class AlertManager {
  private alertRules: AlertRule[] = [];

  constructor() {
    this.setupDefaultAlerts();
  }

  private setupDefaultAlerts(): void {
    // Alertes de disponibilité
    this.addAlertRule({
      name: 'HighErrorRate',
      condition: 'rate(http_requests_total{status=~"5.."}[5m]) > 0.05',
      severity: 'critical',
      duration: 300,
      annotations: {
        summary: 'High error rate detected',
        description: 'Error rate is above 5% for the last 5 minutes'
      },
      labels: { team: 'platform' }
    });

    // Alertes de performance
    this.addAlertRule({
      name: 'HighLatency',
      condition: 'histogram_quantile(0.95, http_request_duration_seconds) > 2',
      severity: 'warning',
      duration: 300,
      annotations: {
        summary: 'High latency detected',
        description: '95th percentile latency is above 2 seconds'
      },
      labels: { team: 'platform' }
    });

    // Alertes spécifiques bancaires
    this.addAlertRule({
      name: 'BankSyncFailure',
      condition: 'rate(bank_sync_operations_total{status="failed"}[10m]) > 0.1',
      severity: 'critical',
      duration: 600,
      annotations: {
        summary: 'Bank synchronization failures',
        description: 'Bank sync failure rate is above 10% for the last 10 minutes'
      },
      labels: { team: 'banking' }
    });

    // Alertes de fraude
    this.addAlertRule({
      name: 'HighFraudDetectionRate',
      condition: 'rate(fraud_detections_total[1h]) > 100',
      severity: 'warning',
      duration: 3600,
      annotations: {
        summary: 'High fraud detection rate',
        description: 'Fraud detection rate is unusually high'
      },
      labels: { team: 'security' }
    });

    // Alertes de business
    this.addAlertRule({
      name: 'LowReconciliationRate',
      condition: 'reconciliation_rate < 0.8',
      severity: 'warning',
      duration: 3600,
      annotations: {
        summary: 'Low reconciliation rate',
        description: 'Reconciliation rate is below 80%'
      },
      labels: { team: 'finance' }
    });
  }

  addAlertRule(rule: AlertRule): void {
    this.alertRules.push(rule);
  }

  getAlertRules(): AlertRule[] {
    return this.alertRules;
  }

  generatePrometheusConfig(): string {
    const groups = this.groupAlertsByTeam();
    let config = 'groups:\n';

    for (const [team, rules] of Object.entries(groups)) {
      config += `  - name: ${team}_alerts\n`;
      config += '    rules:\n';
      
      for (const rule of rules) {
        config += `      - alert: ${rule.name}\n`;
        config += `        expr: ${rule.condition}\n`;
        config += `        for: ${rule.duration}s\n`;
        config += `        labels:\n`;
        
        for (const [key, value] of Object.entries(rule.labels)) {
          config += `          ${key}: "${value}"\n`;
        }
        
        config += `        annotations:\n`;
        for (const [key, value] of Object.entries(rule.annotations)) {
          config += `          ${key}: "${value}"\n`;
        }
      }
    }

    return config;
  }

  private groupAlertsByTeam(): Record<string, AlertRule[]> {
    const grouped: Record<string, AlertRule[]> = {};
    
    for (const rule of this.alertRules) {
      const team = rule.labels.team || 'default';
      if (!grouped[team]) {
        grouped[team] = [];
      }
      grouped[team].push(rule);
    }
    
    return grouped;
  }
}
\end{minted}

\section{Logging Structuré}

\subsection{Format de Log Structuré}

\begin{minted}{typescript}
interface LogEntry {
  timestamp: string;
  level: 'debug' | 'info' | 'warn' | 'error';
  service: string;
  environment: string;
  correlationId?: string;
  userId?: string;
  message: string;
  context?: Record<string, any>;
  error?: {
    name: string;
    message: string;
    stack?: string;
    code?: string;
  };
  tags?: string[];
}

class StructuredLogger {
  constructor(
    private serviceName: string,
    private environment: string
  ) {}

  private createLogEntry(
    level: LogEntry['level'],
    message: string,
    context?: Record<string, any>,
    error?: Error
  ): LogEntry {
    const entry: LogEntry = {
      timestamp: new Date().toISOString(),
      level,
      service: this.serviceName,
      environment: this.environment,
      message
    };

    if (context) {
      entry.context = this.sanitizeContext(context);
    }

    if (error) {
      entry.error = {
        name: error.name,
        message: error.message,
        stack: error.stack,
        code: (error as any).code
      };
    }

    return entry;
  }

  private sanitizeContext(context: Record<string, any>): Record<string, any> {
    const sanitized = { ...context };
    
    // Suppression des données sensibles
    const sensitiveKeys = ['password', 'token', 'secret', 'apiKey', 'creditCard'];
    
    for (const key of Object.keys(sanitized)) {
      if (sensitiveKeys.some(sensitive => key.toLowerCase().includes(sensitive))) {
        sanitized[key] = '[REDACTED]';
      }
    }

    return sanitized;
  }

  info(message: string, context?: Record<string, any>): void {
    const entry = this.createLogEntry('info', message, context);
    console.log(JSON.stringify(entry));
  }

  error(message: string, error?: Error, context?: Record<string, any>): void {
    const entry = this.createLogEntry('error', message, context, error);
    console.error(JSON.stringify(entry));
  }

  warn(message: string, context?: Record<string, any>): void {
    const entry = this.createLogEntry('warn', message, context);
    console.warn(JSON.stringify(entry));
  }

  debug(message: string, context?: Record<string, any>): void {
    if (this.environment === 'development') {
      const entry = this.createLogEntry('debug', message, context);
      console.debug(JSON.stringify(entry));
    }
  }

  // Logging spécifique aux opérations bancaires
  logBankOperation(operation: string, bankId: string, result: any, correlationId: string): void {
    this.info(`Bank operation completed: ${operation}`, {
      bankId,
      operation,
      correlationId,
      result: this.sanitizeBankResult(result)
    });
  }

  private sanitizeBankResult(result: any): any {
    // Suppression des données bancaires sensibles
    const sanitized = { ...result };
    
    if (sanitized.accountNumber) {
      sanitized.accountNumber = this.maskAccountNumber(sanitized.accountNumber);
    }
    
    if (sanitized.transactions) {
      sanitized.transactions = sanitized.transactions.map((tx: any) => ({
        ...tx,
        accountNumber: this.maskAccountNumber(tx.accountNumber)
      }));
    }

    return sanitized;
  }

  private maskAccountNumber(accountNumber: string): string {
    if (!accountNumber || accountNumber.length < 8) return accountNumber;
    return accountNumber.substring(0, 4) + '****' + accountNumber.substring(accountNumber.length - 4);
  }
}
\end{minted}

\subsection{Logging Distribué avec Correlation ID}

\begin{minted}{typescript}
import { v4 as uuidv4 } from 'uuid';

class DistributedLoggingContext {
  private static correlationIdKey = 'x-correlation-id';
  private static userIdKey = 'x-user-id';

  static getCorrelationId(): string {
    // Dans un contexte de requête HTTP, récupérer depuis les headers
    // Pour l'exemple, on génère un nouveau UUID
    return uuidv4();
  }

  static createLoggerMiddleware(logger: StructuredLogger) {
    return (req: any, res: any, next: any) => {
      const correlationId = req.headers[this.correlationIdKey] || uuidv4();
      const userId = req.headers[this.userIdKey];

      // Injection du logger avec contexte dans la requête
      req.logger = {
        info: (message: string, context?: any) => 
          logger.info(message, { ...context, correlationId, userId }),
        error: (message: string, error?: Error, context?: any) => 
          logger.error(message, error, { ...context, correlationId, userId }),
        warn: (message: string, context?: any) => 
          logger.warn(message, { ...context, correlationId, userId })
      };

      // Transmission du correlation ID dans la réponse
      res.setHeader(this.correlationIdKey, correlationId);

      next();
    };
  }
}
\end{minted}

\section{Distributed Tracing}

\subsection{Implémentation avec OpenTelemetry}

\begin{minted}{typescript}
import { NodeTracerProvider } from '@opentelemetry/sdk-trace-node';
import { Resource } from '@opentelemetry/resources';
import { SemanticResourceAttributes } from '@opentelemetry/semantic-conventions';
import { JaegerExporter } from '@opentelemetry/exporter-trace-jaeger';
import { SimpleSpanProcessor } from '@opentelemetry/sdk-trace-base';
import { trace } from '@opentelemetry/api';

class TracingSetup {
  static initialize(serviceName: string, jaegerEndpoint: string): void {
    const resource = Resource.default().merge(
      new Resource({
        [SemanticResourceAttributes.SERVICE_NAME]: serviceName,
        [SemanticResourceAttributes.SERVICE_VERSION]: '1.0.0',
        environment: process.env.NODE_ENV || 'production'
      })
    );

    const provider = new NodeTracerProvider({ resource });

    const exporter = new JaegerExporter({
      endpoint: jaegerEndpoint
    });

    provider.addSpanProcessor(new SimpleSpanProcessor(exporter));
    trace.setGlobalTracerProvider(provider);
  }
}

class TracedBankingService {
  private tracer = trace.getTracer('banking-service');

  async syncBankConnection(connectionId: string): Promise<SyncResult> {
    return this.tracer.startActiveSpan('sync_bank_connection', {
      attributes: {
        'connection.id': connectionId,
        'operation.type': 'sync'
      }
    }, async (span) => {
      try {
        // Span enfant pour l'authentification
        const authResult = await this.tracer.startActiveSpan('bank_authentication', async (authSpan) => {
          authSpan.setAttribute('bank.auth.type', 'oauth2');
          const result = await this.authenticateWithBank(connectionId);
          authSpan.setStatus({ code: 1 }); // OK
          return result;
        });

        // Span enfant pour la récupération des transactions
        const transactions = await this.tracer.startActiveSpan('fetch_transactions', async (fetchSpan) => {
          fetchSpan.setAttribute('fetch.limit', 1000);
          const result = await this.fetchTransactions(connectionId, authResult.token);
          fetchSpan.setAttribute('transaction.count', result.length);
          return result;
        });

        // Span enfant pour le traitement
        const processed = await this.tracer.startActiveSpan('process_transactions', async (processSpan) => {
          const result = await this.processTransactions(transactions);
          processSpan.setAttribute('processed.count', result.length);
          return result;
        });

        span.setStatus({ code: 1 }); // OK
        return { success: true, transactions: processed };

      } catch (error) {
        span.recordException(error as Error);
        span.setStatus({ code: 2, message: (error as Error).message }); // ERROR
        throw error;
      }
    });
  }

  private async authenticateWithBank(connectionId: string): Promise<any> {
    // Implémentation de l'authentification
    return { token: 'sample_token' };
  }

  private async fetchTransactions(connectionId: string, token: string): Promise<any[]> {
    // Implémentation de la récupération
    return [];
  }

  private async processTransactions(transactions: any[]): Promise<any[]> {
    // Implémentation du traitement
    return transactions;
  }
}
\end{minted}

\section{Dashboards et Visualisation}

\subsection{Configuration Grafana}

\begin{minted}{json}
{
  "dashboard": {
    "title": "Financial Platform Overview",
    "panels": [
      {
        "title": "Request Rate",
        "type": "graph",
        "targets": [
          {
            "expr": "rate(http_requests_total[5m])",
            "legendFormat": "{{method}} {{route}}"
          }
        ]
      },
      {
        "title": "Error Rate",
        "type": "graph",
        "targets": [
          {
            "expr": "rate(http_requests_total{status=~\"5..\"}[5m]) / rate(http_requests_total[5m])",
            "legendFormat": "Error Rate"
          }
        ],
        "alert": {
          "conditions": [
            {
              "evaluator": {
                "params": [0.05],
                "type": "gt"
              },
              "operator": {
                "type": "and"
              },
              "query": {
                "params": ["A", "5m", "now"]
              },
              "reducer": {
                "params": [],
                "type": "avg"
              },
              "type": "query"
            }
          ]
        }
      },
      {
        "title": "Bank Sync Operations",
        "type": "stat",
        "targets": [
          {
            "expr": "sum(rate(bank_sync_operations_total[1h]))",
            "legendFormat": "Sync Rate"
          }
        ]
      },
      {
        "title": "Fraud Detections",
        "type": "graph",
        "targets": [
          {
            "expr": "sum by (severity) (rate(fraud_detections_total[1h]))",
            "legendFormat": "{{severity}}"
          }
        ]
      },
      {
        "title": "Transaction Volume",
        "type": "graph",
        "targets": [
          {
            "expr": "sum by (currency) (transaction_volume)",
            "legendFormat": "{{currency}}"
          }
        ]
      },
      {
        "title": "API Latency Distribution",
        "type": "heatmap",
        "targets": [
          {
            "expr": "histogram_quantile(0.95, rate(http_request_duration_seconds_bucket[5m]))",
            "legendFormat": "95th percentile"
          }
        ]
      }
    ]
  }
}
\end{minted}

\begin{bestpractice}
Les dashboards de monitoring doivent être conçus pour répondre à trois questions clés : "Est-ce que le système fonctionne ?", "Y a-t-il des problèmes ?", et "Quelle est la cause du problème ?". Ils doivent permettre une navigation rapide du général au spécifique.
\end{bestpractice}

% ==============================================================================
% CHAPITRE 9 : DIAGNOSTIC ET ANALYSE APPROFONDIE
% ==============================================================================
\chapter{Diagnostic et Analyse Approfondie}

\section{Diagnostic des Ressources et de l'Existant}

\subsection{Analyse de l'Environnement Technique}

L'infrastructure technique de la plateforme repose sur une architecture cloud-native avec les caractéristiques suivantes :

\begin{table}[H]
\centering
\begin{tabularx}{\textwidth}{|X|X|}
\hline
\textbf{Composant} & \textbf{Technologie} \\
\hline
Infrastructure Cloud & AWS / Azure \\
\hline
Orchestration & Kubernetes + Docker \\
\hline
Base de Données & PostgreSQL (relationnelle) + MongoDB (document) \\
\hline
Messagerie & Apache Kafka \\
\hline
Cache & Redis \\
\hline
API Gateway & Kong / AWS API Gateway \\
\hline
Monitoring & Prometheus + Grafana \\
\hline
CI/CD & GitLab CI / GitHub Actions \\
\hline
\end{tabularx}
\caption{Stack Technique de la Plateforme}
\end{table}

\subsection{Diagnostic des Protocoles Bancaires}

\subsubsection{Open Banking (PSD2)}

La directive PSD2 (Payment Services Directive 2) impose l'ouverture des données bancaires via des API sécurisées. Notre implémentation supporte :

\begin{itemize}
    \item \textbf{AIS (Account Information Service)} : Accès aux informations de compte
    \item \textbf{PIS (Payment Initiation Service)} : Initiation de paiements
    \item \textbf{Standardisation} : Conformité aux specifications Berlin Group PSD2
    \item \textbf{Sécurité} : OAuth 2.0 avec PKCE et signatures JWT
\end{itemize}

\subsubsection{EBICS (Electronic Banking Internet Communication Standard)}

EBICS est le standard dominant pour les communications bancaires B2B en Europe. Notre implémentation gère :

\begin{itemize}
    \item \textbf{Protocole sécurisé} : Chiffrement TLS + signature XML
    \item \textbf{Types d'ordres} : HKD, HTD, HAA, CCT, etc.
    \item \textbf{Gestion des clés} : INI, HIA, E002
    \item \textbf{Multi-banques} : Support de plusieurs établissements bancaires
\end{itemize}

\begin{alertbox}
La gestion des clés EBICS nécessite une attention particulière à la sécurité. Les clés privées ne doivent jamais être stockées en clair et doivent être protégées par des HSM (Hardware Security Modules) ou des KMS (Key Management Services).
\end{alertbox}

\section{Architecture Technique du Module F2 -- Multi Banking Hub}

\subsection{Architecture en Couches}

Le module F2 adopte une architecture en couches classique, séparant clairement les responsabilités :

\begin{itemize}
    \item \textbf{Couche API} : Endpoints RESTful avec validation des requêtes
    \item \textbf{Couche Service} : Logique métier et orchestration des protocoles
    \item \textbf{Couche Protocole} : Adaptateurs pour PSD2 et EBICS
    \item \textbf{Couche Data} : Accès aux données avec ORM/ODM
    \item \textbf{Couche Integration} : Communication avec les autres modules via Kafka
\end{itemize}

\subsection{Design Patterns Appliqués}

\subsubsection{Pattern Adapter}

Le pattern Adapter est utilisé pour normaliser les réponses des différents protocoles bancaires :

\begin{minted}{typescript}
interface BankConnectionAdapter {
  connect(credentials: BankCredentials): Promise<ConnectionResult>;
  syncTransactions(connectionId: string): Promise<Transaction[]>;
  getBalance(connectionId: string): Promise<Balance>;
}

class PSD2Adapter implements BankConnectionAdapter {
  async connect(credentials: PSD2Credentials): Promise<ConnectionResult> {
    // Implémentation spécifique PSD2
  }
  
  async syncTransactions(connectionId: string): Promise<Transaction[]> {
    // Synchronisation via API PSD2
  }
}

class EBICSAdapter implements BankConnectionAdapter {
  async connect(credentials: EBICSCredentials): Promise<ConnectionResult> {
    // Implémentation spécifique EBICS
  }
  
  async syncTransactions(connectionId: string): Promise<Transaction[]> {
    // Synchronisation via protocole EBICS
  }
}
\end{minted}

\subsubsection{Pattern Factory}

La Factory Pattern permet de créer dynamiquement le bon adaptateur selon le type de connexion :

\begin{minted}{typescript}
class BankAdapterFactory {
  static createAdapter(type: ConnectionType): BankConnectionAdapter {
    switch (type) {
      case 'PSD2':
        return new PSD2Adapter();
      case 'EBICS':
        return new EBICSAdapter();
      default:
        throw new Error(`Unsupported connection type: ${type}`);
    }
  }
}
\end{minted}

\section{Architecture Technique du Module F5 -- AI Fraud Detection}

\subsection{Architecture Hybride Règles + Machine Learning}

Le module F5 combine une approche basée sur des règles AML/CFT avec des algorithmes de Machine Learning pour une détection de fraude proactive et évolutive.

\subsubsection{Moteur de Règles AML/CFT}

Le moteur de règles implémente les patterns de fraude connus et réglementaires :

\begin{itemize}
    \item \textbf{Règles de montant} : Transactions supérieures à des seuils définis
    \item \textbf{Règles de fréquence} : Nombre de transactions suspectes sur une période
    \item \textbf{Règles géographiques} : Transactions depuis des pays à risque
    \item \textbf{Règles de temps} : Transactions hors heures ouvrables
    \item \textbf{Règles de comportement} : Écarts par rapport au profil historique
\end{itemize}

\subsubsection{Machine Learning Pipeline}

Le pipeline ML comprend plusieurs étapes :

\begin{minted}{python}
from sklearn.ensemble import RandomForestClassifier
from sklearn.preprocessing import StandardScaler
from sklearn.pipeline import Pipeline

class FraudDetectionPipeline:
    def __init__(self):
        self.pipeline = Pipeline([
            ('scaler', StandardScaler()),
            ('classifier', RandomForestClassifier(
                n_estimators=100,
                max_depth=10,
                random_state=42
            ))
        ])
        
    def train(self, X_train, y_train):
        """Entraînement du modèle avec données historiques"""
        self.pipeline.fit(X_train, y_train)
        
    def predict_fraud_probability(self, transaction_features):
        """Prédiction de la probabilité de fraude"""
        return self.pipeline.predict_proba(transaction_features)[:, 1]
        
    def get_feature_importance(self):
        """Importance des features pour l'explicabilité"""
        return self.pipeline.named_steps['classifier'].feature_importances_
\end{minted}

\subsection{Architecture de Graphes Relationnels}

Pour détecter des schémas de fraude complexes (pump-and-dump, money laundering networks), le module utilise une base de données de graphes (Neo4j) :

\begin{minted}{cypher}
// Détection de cycles financiers suspects
MATCH path = (a:Account)-[:TRANSFERRED*]->(a)
WHERE ALL(r IN relationships(path) WHERE r.amount > 10000)
WITH path, length(path) as cycle_length
WHERE cycle_length > 2
RETURN path, cycle_length
ORDER BY cycle_length DESC
LIMIT 100
\end{minted}

\section{Contrats d'Interface OpenAPI}

\subsection{Module F2 -- API Endpoints}

\subsubsection{Création de Connexion Bancaire}

\begin{jsonbox}
\textbf{POST /api/v1/bank-connections}
\end{jsonbox}

\begin{minted}{json}
{
  "correlation_id": "550e8400-e29b-41d4-a716-446655440000",
  "request": {
    "bank_id": "BANK123",
    "connection_type": "PSD2",
    "credentials": {
      "client_id": "client_abc123",
      "client_secret": "secret_xyz789",
      "redirect_uri": "https://app.example.com/callback"
    },
    "scopes": ["ais", "pis"],
    "consent_validity_days": 90
  }
}
\end{minted}

\textbf{Réponse Succès (201 Created)}

\begin{minted}{json}
{
  "correlation_id": "550e8400-e29b-41d4-a716-446655440000",
  "data": {
    "connection_id": "conn_abc123xyz",
    "status": "PENDING_CONSENT",
    "authorization_url": "https://bank.example.com/authorize?client_id=...",
    "expires_at": "2024-03-15T10:30:00Z",
    "created_at": "2024-01-15T10:30:00Z"
  }
}
\end{minted}

\subsubsection{Synchronisation Asynchrone}

\begin{jsonbox}
\textbf{POST /api/v1/bank-connections/\{id\}/sync}
\end{jsonbox}

\begin{minted}{json}
{
  "correlation_id": "660e8400-e29b-41d4-a716-446655440001",
  "sync_options": {
    "from_date": "2024-01-01",
    "to_date": "2024-01-31",
    "include_pending": true,
    "webhook_callback": "https://app.example.com/webhooks/sync"
  }
}
\end{minted}

\textbf{Réponse Succès (202 Accepted)}

\begin{minted}{json}
{
  "correlation_id": "660e8400-e29b-41d4-a716-446655440001",
  "data": {
    "sync_id": "sync_xyz789",
    "status": "PROCESSING",
    "estimated_completion": "2024-01-15T11:00:00Z",
    "message": "Synchronization initiated successfully"
  }
}
\end{minted}

\subsection{Module F5 -- API Endpoints}

\subsubsection{Scan de Fraude}

\begin{jsonbox}
\textbf{POST /api/v1/fraud/scan}
\end{jsonbox}

\begin{minted}{json}
{
  "correlation_id": "770e8400-e29b-41d4-a716-446655440002",
  "scan_config": {
    "transaction_ids": ["txn_001", "txn_002", "txn_003"],
    "scan_types": ["aml_rules", "ml_model", "graph_analysis"],
    "risk_threshold": 0.75,
    "include_explanations": true
  }
}
\end{minted}

\textbf{Réponse Succès (200 OK)}

\begin{minted}{json}
{
  "correlation_id": "770e8400-e29b-41d4-a716-446655440002",
  "data": {
    "scan_id": "scan_fraud_123",
    "total_transactions": 3,
    "flagged_transactions": 1,
    "results": [
      {
        "transaction_id": "txn_002",
        "risk_score": 0.87,
        "risk_level": "HIGH",
        "flagged_rules": ["large_amount", "unusual_location"],
        "ml_probability": 0.92,
        "explanation": "Transaction amount exceeds typical pattern by 450%"
      }
    ],
    "scan_metadata": {
      "executed_at": "2024-01-15T10:35:00Z",
      "processing_time_ms": 245
    }
  }
}
\end{minted}

\subsubsection{Liste des Cas de Fraude}

\begin{jsonbox}
\textbf{GET /api/v1/fraud/cases?status=OPEN\&limit=50}
\end{jsonbox}

\textbf{Réponse Succès (200 OK)}

\begin{minted}{json}
{
  "correlation_id": "880e8400-e29b-41d4-a716-446655440003",
  "data": {
    "total_cases": 15,
    "cases": [
      {
        "case_id": "case_001",
        "status": "OPEN",
        "severity": "HIGH",
        "related_transactions": ["txn_002", "txn_005"]
      }
    ]
  }
}
\end{minted}

\subsection{Gestion Uniforme des Erreurs}

\begin{jsonbox}
\textbf{Format d'Erreur Standard}
\end{jsonbox}

\begin{minted}{json}
{
  "correlation_id": "990e8400-e29b-41d4-a716-446655440004",
  "error": {
    "code": "BANK_CONNECTION_FAILED",
    "message": "Failed to establish connection with bank API",
    "details": {
      "bank_id": "BANK123"
    }
  }
}
\end{minted}

\begin{bestpractice}
L'utilisation d'un \texttt{correlation\_id} unique par requête permet de tracer les appels à travers les différents microservices et de faciliter le debugging en production.
\end{bestpractice}

\section{Analyse de Code et Bonnes Pratiques}

\subsection{Implémentation de la Logique Asynchrone}

La gestion des opérations asynchrones est critique pour les modules F2 et F5, notamment pour les synchronisations bancaires et les scans de fraude qui peuvent prendre plusieurs secondes, voire minutes.

\subsubsection{Pattern Async/Await avec Node.js}

\begin{minted}{typescript}
import { promises as fs } from 'fs';
import { BankConnectionService } from './services';

class BankSyncOrchestrator {
  constructor(private bankService: BankConnectionService) {}

  async syncConnectionWithRetry(
    connectionId: string,
    maxRetries: number = 3
  ): Promise<SyncResult> {
    let attempt = 0;
    const delay = (ms: number) => 
      new Promise(resolve => setTimeout(resolve, ms));

    while (attempt < maxRetries) {
      try {
        console.log(`Attempt ${attempt + 1}`);
        const result = await this.bankService.syncTransactions(connectionId);
        return result;
      } catch (error) {
        attempt++;
        if (attempt === maxRetries) throw error;
        await delay(Math.pow(2, attempt) * 1000); // Backoff exponentiel
      }
    }
  }
}
\end{minted}

% ==============================================================================
% CHAPITRE 10 : CONCLUSION
% ==============================================================================
\chapter{Conclusion}

\section{Synthèse des Réalisations}

Ce rapport a présenté la conception et l'analyse architecturale d'une Plateforme SaaS d'Intelligence Financière globale, avec un focus particulier sur les modules \module{F2 -- Multi Banking Hub} et \module{F5 -- AI Fraud Detection}.

\subsection{Contributions Techniques}

Mon travail a permis de :

\begin{itemize}
    \item \textbf{Concevoir une architecture modulaire} pour le Multi Banking Hub, supportant les protocoles PSD2 et EBICS via des adaptateurs normalisés.
    \item \textbf{Implémenter une détection de fraude hybride} combinant règles AML/CFT, Machine Learning et analyse de graphes relationnels.
    \item \textbf{Définir des contrats d'interface OpenAPI} robustes avec gestion uniforme des erreurs et support de l'asynchronisme.
    \item \textbf{Appliquer les meilleures pratiques} de développement logiciel : idempotence, sécurité OAuth 2.0/JWT, circuit breakers, et monitoring structuré.
\end{itemize}

\subsection{Intégration dans l'Écosystème}

Les modules F2 et F5 jouent un rôle central dans l'écosystème de la plateforme :

\begin{itemize}
    \item Le \module{F2} alimente l'ensemble des modules en données bancaires centralisées et synchronisées.
    \item Le \module{F5} assure la sécurité transversale en validant les transactions et en détectant les schémas de fraude complexes.
\end{itemize}

\section{Perspectives d'Évolution}

Plusieurs pistes d'amélioration pourraient être explorées pour renforcer davantage la plateforme :

\begin{itemize}
    \item \textbf{Extension des protocoles bancaires} : Support de protocoles supplémentaires (SWIFT, Host-to-Host).
    \item \textbf{Amélioration du ML} : Intégration de modèles de deep learning pour la détection de fraude.
    \item \textbf{Edge Computing} : Déploiement de composants en edge computing pour réduire la latence.
\end{itemize}

\section{Conclusion Générale}

Ce projet a permis de mettre en œuvre une architecture logicielle moderne et scalable, répondant aux exigences complexes du secteur financier. L'approche modulaire, couplée à des pratiques de développement rigoureuses, assure la maintenabilité et l'évolutivité de la plateforme.

L'intégration harmonieuse des différents modules démontre l'importance d'une vision architecturale globale dans la conception de systèmes SaaS complexes. La plateforme ainsi conçue constitue une base solide pour l'automatisation financière intelligente et la conformité réglementaire.

\begin{center}
\textit{Ce rapport conclut ma mission de développement et d'analyse architecturale au sein de DATAERA,}
\textit{marquant une contribution significative à la modernisation des solutions financières.}
\end{center}

% ==============================================================================
% ANNEXES
% ==============================================================================
\appendix

\chapter*{Annexes Techniques}
\addcontentsline{toc}{chapter}{Annexes Techniques}

\section{Annexe A : Configuration OpenAPI Complète}

\subsection{Spécification Swagger pour le Module F2}

\begin{minted}{yaml}
openapi: 3.0.0
info:
  title: Multi Banking Hub API
  version: 1.0.0
  description: API pour la gestion des connexions bancaires et la synchronisation des données

servers:
  - url: https://api.example.com/v1
    description: Production server

paths:
  /bank-connections:
    post:
      summary: Créer une nouvelle connexion bancaire
      operationId: createBankConnection
      tags:
        - Bank Connections
      requestBody:
        required: true
        content:
          application/json:
            schema:
              $ref: '#/components/schemas/CreateConnectionRequest'
      responses:
        '201':
          description: Connexion créée avec succès
          content:
            application/json:
              schema:
                $ref: '#/components/schemas/ConnectionResponse'
        '400':
          description: Requête invalide
          content:
            application/json:
              schema:
                $ref: '#/components/schemas/Error'
        '401':
          description: Non autorisé
          content:
            application/json:
              schema:
                $ref: '#/components/schemas/Error'

  /bank-connections/{connectionId}/sync:
    post:
      summary: Synchroniser les données bancaires
      operationId: syncBankData
      tags:
        - Bank Connections
      parameters:
        - name: connectionId
          in: path
          required: true
          schema:
            type: string
      requestBody:
        required: true
        content:
          application/json:
            schema:
              $ref: '#/components/schemas/SyncRequest'
      responses:
        '202':
          description: Synchronisation acceptée
          content:
            application/json:
              schema:
                $ref: '#/components/schemas/SyncResponse'
        '404':
          description: Connexion non trouvée
          content:
            application/json:
              schema:
                $ref: '#/components/schemas/Error'

components:
  schemas:
    CreateConnectionRequest:
      type: object
      required:
        - bank_id
        - connection_type
        - credentials
      properties:
        correlation_id:
          type: string
          format: uuid
        bank_id:
          type: string
        connection_type:
          type: string
          enum: [PSD2, EBICS]
        credentials:
          type: object
        scopes:
          type: array
          items:
            type: string
        consent_validity_days:
          type: integer
          minimum: 1
          maximum: 365

    ConnectionResponse:
      type: object
      properties:
        correlation_id:
          type: string
        data:
          type: object
          properties:
            connection_id:
              type: string
            status:
              type: string
            authorization_url:
              type: string
            expires_at:
              type: string
              format: date-time
            created_at:
              type: string
              format: date-time

    SyncRequest:
      type: object
      properties:
        correlation_id:
          type: string
        sync_options:
          type: object
          properties:
            from_date:
              type: string
              format: date
            to_date:
              type: string
              format: date
            include_pending:
              type: boolean
            webhook_callback:
              type: string
              format: uri

    Error:
      type: object
      properties:
        correlation_id:
          type: string
        error:
          type: object
          properties:
            code:
              type: string
            message:
              type: string
            details:
              type: object
\end{minted}

\section{Annexe B : Scripts de Déploiement Kubernetes}

\subsection{Deployment YAML pour le Module F2}

\begin{minted}{yaml}
apiVersion: apps/v1
kind: Deployment
metadata:
  name: bank-connection-service
  labels:
    app: bank-connection
    version: v1
spec:
  replicas: 3
  selector:
    matchLabels:
      app: bank-connection
  template:
    metadata:
      labels:
        app: bank-connection
        version: v1
    spec:
      containers:
      - name: bank-connection
        image: registry.example.com/bank-connection:1.0.0
        ports:
        - containerPort: 8080
          name: http
        env:
        - name: DATABASE_URL
          valueFrom:
            secretKeyRef:
              name: database-secrets
              key: url
        - name: REDIS_URL
          valueFrom:
            secretKeyRef:
              name: redis-secrets
              key: url
        - name: KAFKA_BROKERS
          value: "kafka:9092"
        - name: JWT_ISSUER
          value: "https://auth.example.com"
        resources:
          requests:
            memory: "256Mi"
            cpu: "250m"
          limits:
            memory: "512Mi"
            cpu: "500m"
        livenessProbe:
          httpGet:
            path: /health
            port: 8080
          initialDelaySeconds: 30
          periodSeconds: 10
        readinessProbe:
          httpGet:
            path: /ready
            port: 8080
          initialDelaySeconds: 5
          periodSeconds: 5
---
apiVersion: v1
kind: Service
metadata:
  name: bank-connection-service
spec:
  selector:
    app: bank-connection
  ports:
  - protocol: TCP
    port: 80
    targetPort: 8080
  type: ClusterIP
---
apiVersion: autoscaling/v2
kind: HorizontalPodAutoscaler
metadata:
  name: bank-connection-hpa
spec:
  scaleTargetRef:
    apiVersion: apps/v1
    kind: Deployment
    name: bank-connection-service
  minReplicas: 2
  maxReplicas: 10
  metrics:
  - type: Resource
    resource:
      name: cpu
      target:
        type: Utilization
        averageUtilization: 70
  - type: Resource
    resource:
      name: memory
      target:
        type: Utilization
        averageUtilization: 80
\end{minted}

\section{Annexe C : Configuration CI/CD}

\subsection{Pipeline GitLab CI}

\begin{minted}{yaml}
stages:
  - test
  - build
  - deploy

variables:
  DOCKER_IMAGE: registry.example.com/bank-connection
  DOCKER_TAG: $CI_COMMIT_SHORT_SHA

test:
  stage: test
  image: node:18
  script:
    - npm ci
    - npm run lint
    - npm run test:unit
    - npm run test:integration
  coverage: '/All files[^|]*\|[^|]*\s+([\d\.]+)/'
  artifacts:
    reports:
      coverage_report:
        coverage_format: cobertura
        path: coverage/cobertura-coverage.xml

build:
  stage: build
  image: docker:24
  services:
    - docker:24-dind
  script:
    - docker login -u $CI_REGISTRY_USER -p $CI_REGISTRY_PASSWORD $CI_REGISTRY
    - docker build -t $DOCKER_IMAGE:$DOCKER_TAG .
    - docker push $DOCKER_IMAGE:$DOCKER_TAG
    - docker tag $DOCKER_IMAGE:$DOCKER_TAG $DOCKER_IMAGE:latest
    - docker push $DOCKER_IMAGE:latest
  only:
    - main
    - develop

deploy:staging:
  stage: deploy
  image: bitnami/kubectl:latest
  script:
    - kubectl config use-context staging
    - kubectl set image deployment/bank-connection-service bank-connection=$DOCKER_IMAGE:$DOCKER_TAG
    - kubectl rollout status deployment/bank-connection-service
  environment:
    name: staging
    url: https://staging.example.com
  only:
    - develop

deploy:production:
  stage: deploy
  image: bitnami/kubectl:latest
  script:
    - kubectl config use-context production
    - kubectl set image deployment/bank-connection-service bank-connection=$DOCKER_IMAGE:$DOCKER_TAG
    - kubectl rollout status deployment/bank-connection-service
  environment:
    name: production
    url: https://api.example.com
  when: manual
  only:
    - main
\end{minted}

\section{Annexe D : Glossaire Technique}

\begin{table}[H]
\centering
\begin{tabularx}{\textwidth}{|X|X|}
\hline
\textbf{Terme} & \textbf{Définition} \\
\hline
AIS & Account Information Service - Service d'information de compte PSD2 \\
\hline
AML & Anti-Money Laundering - Lutte contre le blanchiment d'argent \\
\hline
API & Application Programming Interface - Interface de programmation \\
\hline
CFT & Combating the Financing of Terrorism - Lutte contre le financement du terrorisme \\
\hline
EBICS & Electronic Banking Internet Communication Standard \\
\hline
GDPR & General Data Protection Regulation - Règlement sur la protection des données \\
\hline
HSM & Hardware Security Module - Module de sécurité matériel \\
\hline
JWT & JSON Web Token - Token Web JSON \\
\hline
KMS & Key Management Service - Service de gestion des clés \\
\hline
OAuth 2.0 & Open Authorization 2.0 - Protocole d'autorisation \\
\hline
PCI DSS & Payment Card Industry Data Security Standard \\
\hline
PIS & Payment Initiation Service - Service d'initiation de paiement PSD2 \\
\hline
PKCE & Proof Key for Code Exchange - Preuve de clé par échange de code \\
\hline
PSD2 & Payment Services Directive 2 - Directive sur les services de paiement \\
\hline
RAG & Retrieval-Augmented Generation - Génération augmentée par récupération \\
\hline
RBAC & Role-Based Access Control - Contrôle d'accès basé sur les rôles \\
\hline
REST & Representational State Transfer - Transfert d'état représentationnel \\
\hline
SaaS & Software as a Service - Logiciel en tant que service \\
\hline
SCA & Strong Customer Authentication - Authentification forte du client \\
\hline
SHAP & SHapley Additive exPlanations - Explications additives SHapley \\
\hline
SOC 2 & Service Organization Control 2 - Contrôle d'organisation de service \\
\hline
\end{tabularx}
\caption{Glossaire Technique}
\end{table}

% ==============================================================================
% BIBLIOGRAPHIE
% ==============================================================================
\chapter*{Bibliographie}
\addcontentsline{toc}{chapter}{Bibliographie}

\begin{enumerate}
    \item European Banking Authority. (2019). \textit{Regulatory Technical Standards on Strong Customer Authentication and Common Secure Communication}.
    \item Berlin Group. (2023). \textit{NextGenPSD2 Framework Specification}.
    \item Deutsche Kreditwirtschaft. (2022). \textit{EBICS Specification Version 3.0}.
    \item Financial Action Task Force (FATF). (2023). \textit{International Standards on Combating Money Laundering and the Financing of Terrorism}.
    \item Newman, S. (2021). \textit{Building Microservices: Designing Fine-Grained Systems}. O'Reilly Media.
    \item Richardson, L., & Amundsen, M. (2022). \textit{RESTful Web APIs}. O'Reilly Media.
    \item Gamma, E., Helm, R., Johnson, R., & Vlissides, J. (1994). \textit{Design Patterns: Elements of Reusable Object-Oriented Software}. Addison-Wesley.
    \item Chollet, F. (2021). \textit{Deep Learning with Python}. Manning Publications.
    \item Molnar, C. (2022). \textit{Interpretable Machine Learning}. Lulu.com.
    \item ISO/IEC 27001. (2022). \textit{Information security, cybersecurity and privacy protection — Information security management systems}.
\end{enumerate}
\end{minted}

\begin{bestpractice}
Le monitoring de sécurité en temps réel permet de détecter et de répondre aux menaces avant qu'elles ne causent des dommages significatifs. Les règles de détection doivent être régulièrement ajustées basées sur les patterns observés.
\end{bestpractice}
\end{minted}

% ==============================================================================
% CHAPITRE 3 : DIAGNOSTIC ET ANALYSE APPROFONDIE
% ==============================================================================
\chapter{Diagnostic et Analyse Approfondie}

\section{Diagnostic des Ressources et de l'Existant}

\subsection{Analyse de l'Environnement Technique}

L'infrastructure technique de la plateforme repose sur une architecture cloud-native avec les caractéristiques suivantes :

\begin{table}[H]
\centering
\begin{tabularx}{\textwidth}{|X|X|}
\hline
\textbf{Composant} & \textbf{Technologie} \\
\hline
Infrastructure Cloud & AWS / Azure \\
\hline
Orchestration & Kubernetes + Docker \\
\hline
Base de Données & PostgreSQL (relationnelle) + MongoDB (document) \\
\hline
Messagerie & Apache Kafka \\
\hline
Cache & Redis \\
\hline
API Gateway & Kong / AWS API Gateway \\
\hline
Monitoring & Prometheus + Grafana \\
\hline
CI/CD & GitLab CI / GitHub Actions \\
\hline
\end{tabularx}
\caption{Stack Technique de la Plateforme}
\end{table}

\subsection{Diagnostic des Protocoles Bancaires}

\subsubsection{Open Banking (PSD2)}

La directive PSD2 (Payment Services Directive 2) impose l'ouverture des données bancaires via des API sécurisées. Notre implémentation supporte :

\begin{itemize}
    \item \textbf{AIS (Account Information Service)} : Accès aux informations de compte
    \item \textbf{PIS (Payment Initiation Service)} : Initiation de paiements
    \item \textbf{Standardisation} : Conformité aux specifications Berlin Group PSD2
    \item \textbf{Sécurité} : OAuth 2.0 avec PKCE et signatures JWT
\end{itemize}

\subsubsection{EBICS (Electronic Banking Internet Communication Standard)}

EBICS est le standard dominant pour les communications bancaires B2B en Europe. Notre implémentation gère :

\begin{itemize}
    \item \textbf{Protocole sécurisé} : Chiffrement TLS + signature XML
    \item \textbf{Types d'ordres} : HKD, HTD, HAA, CCT, etc.
    \item \textbf{Gestion des clés} : INI, HIA, E002
    \item \textbf{Multi-banques} : Support de plusieurs établissements bancaires
\end{itemize}

\begin{alertbox}
La gestion des clés EBICS nécessite une attention particulière à la sécurité. Les clés privées ne doivent jamais être stockées en clair et doivent être protégées par des HSM (Hardware Security Modules) ou des KMS (Key Management Services).
\end{alertbox}

\section{Architecture Technique du Module F2 -- Multi Banking Hub}

\subsection{Architecture en Couches}

Le module F2 adopte une architecture en couches classique, séparant clairement les responsabilités :

\begin{itemize}
    \item \textbf{Couche API} : Endpoints RESTful avec validation des requêtes
    \item \textbf{Couche Service} : Logique métier et orchestration des protocoles
    \item \textbf{Couche Protocole} : Adaptateurs pour PSD2 et EBICS
    \item \textbf{Couche Data} : Accès aux données avec ORM/ODM
    \item \textbf{Couche Integration} : Communication avec les autres modules via Kafka
\end{itemize}

\subsection{Design Patterns Appliqués}

\subsubsection{Pattern Adapter}

Le pattern Adapter est utilisé pour normaliser les réponses des différents protocoles bancaires :

\begin{minted}{typescript}
interface BankConnectionAdapter {
  connect(credentials: BankCredentials): Promise<ConnectionResult>;
  syncTransactions(connectionId: string): Promise<Transaction[]>;
  getBalance(connectionId: string): Promise<Balance>;
}

class PSD2Adapter implements BankConnectionAdapter {
  async connect(credentials: PSD2Credentials): Promise<ConnectionResult> {
    // Implémentation spécifique PSD2
  }
  
  async syncTransactions(connectionId: string): Promise<Transaction[]> {
    // Synchronisation via API PSD2
  }
}

class EBICSAdapter implements BankConnectionAdapter {
  async connect(credentials: EBICSCredentials): Promise<ConnectionResult> {
    // Implémentation spécifique EBICS
  }
  
  async syncTransactions(connectionId: string): Promise<Transaction[]> {
    // Synchronisation via protocole EBICS
  }
}
\end{minted}

\subsubsection{Pattern Factory}

La Factory Pattern permet de créer dynamiquement le bon adaptateur selon le type de connexion :

\begin{minted}{typescript}
class BankAdapterFactory {
  static createAdapter(type: ConnectionType): BankConnectionAdapter {
    switch (type) {
      case 'PSD2':
        return new PSD2Adapter();
      case 'EBICS':
        return new EBICSAdapter();
      default:
        throw new Error(`Unsupported connection type: ${type}`);
    }
  }
}
\end{minted}

\section{Architecture Technique du Module F5 -- AI Fraud Detection}

\subsection{Architecture Hybride Règles + Machine Learning}

Le module F5 combine une approche basée sur des règles AML/CFT avec des algorithmes de Machine Learning pour une détection de fraude proactive et évolutive.

\subsubsection{Moteur de Règles AML/CFT}

Le moteur de règles implémente les patterns de fraude connus et réglementaires :

\begin{itemize}
    \item \textbf{Règles de montant} : Transactions supérieures à des seuils définis
    \item \textbf{Règles de fréquence} : Nombre de transactions suspectes sur une période
    \item \textbf{Règles géographiques} : Transactions depuis des pays à risque
    \item \textbf{Règles de temps} : Transactions hors heures ouvrables
    \item \textbf{Règles de comportement} : Écarts par rapport au profil historique
\end{itemize}

\subsubsection{Machine Learning Pipeline}

Le pipeline ML comprend plusieurs étapes :

\begin{minted}{python}
from sklearn.ensemble import RandomForestClassifier
from sklearn.preprocessing import StandardScaler
from sklearn.pipeline import Pipeline

class FraudDetectionPipeline:
    def __init__(self):
        self.pipeline = Pipeline([
            ('scaler', StandardScaler()),
            ('classifier', RandomForestClassifier(
                n_estimators=100,
                max_depth=10,
                random_state=42
            ))
        ])
        
    def train(self, X_train, y_train):
        """Entraînement du modèle avec données historiques"""
        self.pipeline.fit(X_train, y_train)
        
    def predict_fraud_probability(self, transaction_features):
        """Prédiction de la probabilité de fraude"""
        return self.pipeline.predict_proba(transaction_features)[:, 1]
        
    def get_feature_importance(self):
        """Importance des features pour l'explicabilité"""
        return self.pipeline.named_steps['classifier'].feature_importances_
\end{minted}

\subsection{Architecture de Graphes Relationnels}

Pour détecter des schémas de fraude complexes (pump-and-dump, money laundering networks), le module utilise une base de données de graphes (Neo4j) :

\begin{minted}{cypher}
// Détection de cycles financiers suspects
MATCH path = (a:Account)-[:TRANSFERRED*]->(a)
WHERE ALL(r IN relationships(path) WHERE r.amount > 10000)
WITH path, length(path) as cycle_length
WHERE cycle_length > 2
RETURN path, cycle_length
ORDER BY cycle_length DESC
LIMIT 100
\end{minted}

\section{Contrats d'Interface OpenAPI}

\subsection{Module F2 -- API Endpoints}

\subsubsection{Création de Connexion Bancaire}

\begin{jsonbox}
\textbf{POST /api/v1/bank-connections}
\end{jsonbox}

\begin{minted}{json}
{
  "correlation_id": "550e8400-e29b-41d4-a716-446655440000",
  "request": {
    "bank_id": "BANK123",
    "connection_type": "PSD2",
    "credentials": {
      "client_id": "client_abc123",
      "client_secret": "secret_xyz789",
      "redirect_uri": "https://app.example.com/callback"
    },
    "scopes": ["ais", "pis"],
    "consent_validity_days": 90
  }
}
\end{minted}

\textbf{Réponse Succès (201 Created)}

\begin{minted}{json}
{
  "correlation_id": "550e8400-e29b-41d4-a716-446655440000",
  "data": {
    "connection_id": "conn_abc123xyz",
    "status": "PENDING_CONSENT",
    "authorization_url": "https://bank.example.com/authorize?client_id=...",
    "expires_at": "2024-03-15T10:30:00Z",
    "created_at": "2024-01-15T10:30:00Z"
  }
}
\end{minted}

\subsubsection{Synchronisation Asynchrone}

\begin{jsonbox}
\textbf{POST /api/v1/bank-connections/\{id\}/sync}
\end{jsonbox}

\begin{minted}{json}
{
  "correlation_id": "660e8400-e29b-41d4-a716-446655440001",
  "sync_options": {
    "from_date": "2024-01-01",
    "to_date": "2024-01-31",
    "include_pending": true,
    "webhook_callback": "https://app.example.com/webhooks/sync"
  }
}
\end{minted}

\textbf{Réponse Succès (202 Accepted)}

\begin{minted}{json}
{
  "correlation_id": "660e8400-e29b-41d4-a716-446655440001",
  "data": {
    "sync_id": "sync_xyz789",
    "status": "PROCESSING",
    "estimated_completion": "2024-01-15T11:00:00Z",
    "message": "Synchronization initiated successfully"
  }
}
\end{minted}

\subsection{Module F5 -- API Endpoints}

\subsubsection{Scan de Fraude}

\begin{jsonbox}
\textbf{POST /api/v1/fraud/scan}
\end{jsonbox}

\begin{minted}{json}
{
  "correlation_id": "770e8400-e29b-41d4-a716-446655440002",
  "scan_config": {
    "transaction_ids": ["txn_001", "txn_002", "txn_003"],
    "scan_types": ["aml_rules", "ml_model", "graph_analysis"],
    "risk_threshold": 0.75,
    "include_explanations": true
  }
}
\end{minted}

\textbf{Réponse Succès (200 OK)}

\begin{minted}{json}
{
  "correlation_id": "770e8400-e29b-41d4-a716-446655440002",
  "data": {
    "scan_id": "scan_fraud_123",
    "total_transactions": 3,
    "flagged_transactions": 1,
    "results": [
      {
        "transaction_id": "txn_002",
        "risk_score": 0.87,
        "risk_level": "HIGH",
        "flagged_rules": ["large_amount", "unusual_location"],
        "ml_probability": 0.92,
        "explanation": "Transaction amount exceeds typical pattern by 450%"
      }
    ],
    "scan_metadata": {
      "executed_at": "2024-01-15T10:35:00Z",
      "processing_time_ms": 245
    }
  }
}
\end{minted}

\subsubsection{Liste des Cas de Fraude}

\begin{jsonbox}
\textbf{GET /api/v1/fraud/cases?status=OPEN\&limit=50}
\end{jsonbox}

\textbf{Réponse Succès (200 OK)}

\begin{minted}{json}
{
  "correlation_id": "880e8400-e29b-41d4-a716-446655440003",
  "data": {
    "total_cases": 15,
    "cases": [
      {
        "case_id": "case_001",
        "status": "OPEN",
        "severity": "HIGH",
        "created_at": "2024-01-14T14:30:00Z",
        "related_transactions": ["txn_002", "txn_005"],
        "assigned_to": "analyst_1",
        "notes": "Potential money laundering pattern detected"
      }
    ],
    "pagination": {
      "page": 1,
      "limit": 50,
      "total_pages": 1
    }
  }
}
\end{minted}

\subsection{Gestion Uniforme des Erreurs}

Tous les endpoints retournent un format d'erreur standardisé pour faciliter le debugging et le monitoring :

\begin{jsonbox}
\textbf{Format d'Erreur Standard}
\end{jsonbox}

\begin{minted}{json}
{
  "correlation_id": "990e8400-e29b-41d4-a716-446655440004",
  "error": {
    "code": "BANK_CONNECTION_FAILED",
    "message": "Failed to establish connection with bank API",
    "details": {
      "bank_id": "BANK123",
      "error_code": "AUTH_INVALID",
      "retryable": true,
      "retry_after_seconds": 60
    },
    "timestamp": "2024-01-15T10:40:00Z",
    "request_id": "req_abc123"
  }
}
\end{minted}

\begin{bestpractice}
L'utilisation d'un \texttt{correlation\_id} unique par requête permet de tracer les appels à travers les différents microservices et de faciliter le debugging en production.
\end{bestpractice}

\section{Analyse de Code et Bonnes Pratiques}

\subsection{Implémentation de la Logique Asynchrone}

La gestion des opérations asynchrones est critique pour les modules F2 et F5, notamment pour les synchronisations bancaires et les scans de fraude qui peuvent prendre plusieurs secondes, voire minutes.

\subsubsection{Pattern Async/Await avec Node.js}

\begin{minted}{typescript}
import { promises as fs } from 'fs';
import { BankConnectionService } from './services';

class BankSyncOrchestrator {
  constructor(private bankService: BankConnectionService) {}

  async syncConnectionWithRetry(
    connectionId: string,
    maxRetries: number = 3
  ): Promise<SyncResult> {
    let attempt = 0;
    const delay = (ms: number) => 
      new Promise(resolve => setTimeout(resolve, ms));

    while (attempt < maxRetries) {
      try {
        console.log(`Attempt ${attempt + 1} for connection ${connectionId}`);
        
        const result = await this.bankService.syncTransactions(connectionId);
        
        // Publication de l'événement de succès
        await this.publishSyncEvent({
          connectionId,
          status: 'SUCCESS',
          transactionCount: result.transactions.length
        });
        
        return result;
        
      } catch (error) {
        attempt++;
        
        if (attempt === maxRetries) {
          // Publication de l'événement d'échec
          await this.publishSyncEvent({
            connectionId,
            status: 'FAILED',
            error: error.message
          });
          throw new Error(
            `Sync failed after ${maxRetries} attempts: ${error.message}`
          );
        }
        
        // Backoff exponentiel
        const backoffTime = Math.pow(2, attempt) * 1000;
        console.log(`Retrying in ${backoffTime}ms...`);
        await delay(backoffTime);
      }
    }
  }

  private async publishSyncEvent(event: SyncEvent): Promise<void> {
    // Implémentation de la publication Kafka
  }
}
\end{minted}

\subsubsection{Pattern Circuit Breaker}

Pour éviter les cascades d'échecs, nous implémentons un Circuit Breaker :

\begin{minted}{typescript}
enum CircuitState {
  CLOSED = 'CLOSED',
  OPEN = 'OPEN',
  HALF_OPEN = 'HALF_OPEN'
}

class CircuitBreaker {
  private state: CircuitState = CircuitState.CLOSED;
  private failureCount = 0;
  private lastFailureTime?: Date;
  private successCount = 0;

  constructor(
    private threshold: number = 5,
    private timeout: number = 60000 // 1 minute
  ) {}

  async execute<T>(
    operation: () => Promise<T>
  ): Promise<T> {
    if (this.state === CircuitState.OPEN) {
      if (this.shouldAttemptReset()) {
        this.state = CircuitState.HALF_OPEN;
      } else {
        throw new Error('Circuit breaker is OPEN');
      }
    }

    try {
      const result = await operation();
      this.onSuccess();
      return result;
    } catch (error) {
      this.onFailure();
      throw error;
    }
  }

  private onSuccess(): void {
    this.failureCount = 0;
    if (this.state === CircuitState.HALF_OPEN) {
      this.successCount++;
      if (this.successCount >= 2) {
        this.state = CircuitState.CLOSED;
        this.successCount = 0;
      }
    }
  }

  private onFailure(): void {
    this.failureCount++;
    this.lastFailureTime = new Date();
    
    if (this.failureCount >= this.threshold) {
      this.state = CircuitState.OPEN;
    }
  }

  private shouldAttemptReset(): boolean {
    return this.lastFailureTime && 
      Date.now() - this.lastFailureTime.getTime() > this.timeout;
  }
}
\end{minted}

\subsection{Implémentation de l'Idempotence}

L'idempotence est cruciale pour éviter les traitements en double lors des reconnexions réseau ou des retentatives automatiques.

\subsubsection{Header Idempotency-Key}

\begin{minted}{typescript}
import { createHash } from 'crypto';

class IdempotencyManager {
  private cache: Map<string, CachedResponse> = new Map();
  private readonly TTL = 24 * 60 * 60 * 1000; // 24 heures

  async processRequest<T>(
    idempotencyKey: string,
    operation: () => Promise<T>
  ): Promise<T> {
    // Vérification du cache
    const cached = this.cache.get(idempotencyKey);
    if (cached && !this.isExpired(cached)) {
      console.log(`Returning cached response for key: ${idempotencyKey}`);
      return cached.response as T;
    }

    // Exécution de l'opération
    const result = await operation();

    // Mise en cache
    this.cache.set(idempotencyKey, {
      response: result,
      timestamp: Date.now()
    });

    return result;
  }

  private isExpired(cached: CachedResponse): boolean {
    return Date.now() - cached.timestamp > this.TTL;
  }

  generateIdempotencyKey(payload: any): string {
    const hash = createHash('sha256');
    hash.update(JSON.stringify(payload));
    return hash.digest('hex');
  }
}

// Middleware Express
export function idempotencyMiddleware(
  manager: IdempotencyManager
) {
  return async (req: any, res: any, next: any) => {
    const idempotencyKey = req.headers['idempotency-key'];
    
    if (!idempotencyKey) {
      return next();
    }

    try {
      const cached = await manager.processRequest(
        idempotencyKey,
        () => new Promise((resolve) => {
          // Stockage de la réponse originale
          const originalJson = res.json.bind(res);
          res.json = (data: any) => {
            resolve(data);
            return originalJson(data);
          };
          next();
        })
      );
      
      if (cached) {
        return res.status(200).json(cached);
      }
    } catch (error) {
      next(error);
    }
  };
}
\end{minted}

\begin{bestpractice}
L'implémentation de l'idempotence via le header \texttt{Idempotency-Key} permet aux clients de réitérer en toute sécurité les mêmes requêtes sans risque de duplication des traitements.
\end{bestpractice}

\subsection{Implémentation de la Sécurité OAuth 2.0 / JWT}

La sécurité des API est primordiale pour une plateforme financière. Nous utilisons OAuth 2.0 avec JWT pour l'authentification et l'autorisation.

\subsubsection{Validation des Tokens JWT}

\begin{minted}{typescript}
import jwt from 'jsonwebtoken';
import jwksClient from 'jwks-rsa';

class JWTValidator {
  private client: any;

  constructor(jwksUri: string) {
    this.client = jwksClient({
      jwksUri: jwksUri,
      cache: true,
      rateLimit: true,
      jwksRequestsPerMinute: 10
    });
  }

  async validateToken(token: string): Promise<DecodedToken> {
    try {
      const decoded = jwt.decode(token, { complete: true });
      
      if (!decoded || !decoded.header) {
        throw new Error('Invalid token format');
      }

      const key = await this.getSigningKey(decoded.header.kid);
      
      const verified = jwt.verify(token, key, {
        algorithms: ['RS256'],
        issuer: process.env.JWT_ISSUER,
        audience: process.env.JWT_AUDIENCE
      });

      return verified as DecodedToken;
      
    } catch (error) {
      throw new Error(`Token validation failed: ${error.message}`);
    }
  }

  private async getSigningKey(kid: string): Promise<string> {
    return new Promise((resolve, reject) => {
      this.client.getSigningKey(kid, (err: any, key: any) => {
        if (err) {
          reject(err);
        } else {
          resolve(key.publicKey || key.rsaPublicKey);
        }
      });
    });
  }
}

// Middleware d'authentification
export function authMiddleware(validator: JWTValidator) {
  return async (req: any, res: any, next: any) => {
    const authHeader = req.headers.authorization;
    
    if (!authHeader || !authHeader.startsWith('Bearer ')) {
      return res.status(401).json({
        error: 'Unauthorized',
        message: 'Missing or invalid authorization header'
      });
    }

    const token = authHeader.substring(7);
    
    try {
      const decoded = await validator.validateToken(token);
      
      // Ajout des informations utilisateur à la requête
      req.user = {
        id: decoded.sub,
        email: decoded.email,
        roles: decoded.roles,
        permissions: decoded.permissions
      };
      
      next();
    } catch (error) {
      return res.status(401).json({
        error: 'Unauthorized',
        message: error.message
      });
    }
  };
}
\end{minted}

\subsubsection{Contrôle d'Accès Basé sur les Rôles (RBAC)}

\begin{minted}{typescript}
enum Permission {
  READ_TRANSACTIONS = 'read:transactions',
  WRITE_TRANSACTIONS = 'write:transactions',
  MANAGE_CONNECTIONS = 'manage:connections',
  VIEW_FRAUD_CASES = 'view:fraud_cases',
  RESOLVE_FRAUD_CASES = 'resolve:fraud_cases'
}

class AuthorizationManager {
  constructor(private requiredPermissions: Permission[]) {}

  authorize(user: any): boolean {
    if (!user.permissions || !Array.isArray(user.permissions)) {
      return false;
    }

    return this.requiredPermissions.every(permission =>
      user.permissions.includes(permission)
    );
  }
}

// Middleware d'autorisation
export function requirePermissions(...permissions: Permission[]) {
  const authzManager = new AuthorizationManager(permissions);
  
  return (req: any, res: any, next: any) => {
    if (!req.user) {
      return res.status(401).json({
        error: 'Unauthorized',
        message: 'Authentication required'
      });
    }

    if (!authzManager.authorize(req.user)) {
      return res.status(403).json({
        error: 'Forbidden',
        message: 'Insufficient permissions'
      });
    }

    next();
  };
}

// Utilisation dans les routes
app.post(
  '/api/v1/bank-connections',
  authMiddleware(jwtValidator),
  requirePermissions(Permission.MANAGE_CONNECTIONS),
  createBankConnectionHandler
);
\end{minted}

\begin{alertbox}
Les tokens JWT doivent avoir une durée de vie limitée (15-30 minutes) et être rafraîchis via des refresh tokens stockés de manière sécurisée (HttpOnly cookies).
\end{alertbox}

\subsection{Gestion des Secrets et Configuration}

\subsubsection{Utilisation de Variables d'Environnement}

\begin{minted}{typescript}
// config.ts
interface AppConfig {
  database: {
    host: string;
    port: number;
    username: string;
    password: string;
    database: string;
  };
  redis: {
    host: string;
    port: number;
    password: string;
  };
  jwt: {
    issuer: string;
    audience: string;
    jwksUri: string;
  };
  banking: {
    psd2: {
      clientId: string;
      clientSecret: string;
    };
    ebics: {
      host: string;
      partnerId: string;
      userId: string;
    };
  };
}

export function loadConfig(): AppConfig {
  return {
    database: {
      host: process.env.DB_HOST || 'localhost',
      port: parseInt(process.env.DB_PORT || '5432'),
      username: process.env.DB_USERNAME,
      password: process.env.DB_PASSWORD,
      database: process.env.DB_NAME
    },
    redis: {
      host: process.env.REDIS_HOST || 'localhost',
      port: parseInt(process.env.REDIS_PORT || '6379'),
      password: process.env.REDIS_PASSWORD
    },
    jwt: {
      issuer: process.env.JWT_ISSUER,
      audience: process.env.JWT_AUDIENCE,
      jwksUri: process.env.JWT_JWKS_URI
    },
    banking: {
      psd2: {
        clientId: process.env.PSD2_CLIENT_ID,
        clientSecret: process.env.PSD2_CLIENT_SECRET
      },
      ebics: {
        host: process.env.EBICS_HOST,
        partnerId: process.env.EBICS_PARTNER_ID,
        userId: process.env.EBICS_USER_ID
      }
    }
  };
}

// Validation de la configuration
export function validateConfig(config: AppConfig): void {
  const requiredFields = [
    'database.username',
    'database.password',
    'jwt.issuer',
    'banking.psd2.clientId'
  ];

  for (const field of requiredFields) {
    const value = field.split('.').reduce((obj, key) => obj?.[key], config);
    if (!value) {
      throw new Error(`Missing required configuration: ${field}`);
    }
  }
}
\end{minted}

\subsection{Monitoring et Logging}

\subsubsection{Structured Logging}

\begin{minted}{typescript}
interface LogContext {
  correlationId?: string;
  userId?: string;
  service: string;
  module: string;
  action: string;
}

class Logger {
  constructor(private service: string) {}

  private formatMessage(
    level: string,
    message: string,
    context?: LogContext,
    data?: any
  ): string {
    return JSON.stringify({
      timestamp: new Date().toISOString(),
      level,
      service: this.service,
      message,
      ...context,
      ...(data && { data })
    });
  }

  info(message: string, context?: LogContext, data?: any): void {
    console.log(this.formatMessage('INFO', message, context, data));
  }

  error(message: string, context?: LogContext, error?: any): void {
    console.error(this.formatMessage('ERROR', message, context, {
      error: error?.message,
      stack: error?.stack
    }));
  }

  warn(message: string, context?: LogContext, data?: any): void {
    console.warn(this.formatMessage('WARN', message, context, data));
  }
}

// Utilisation
const logger = new Logger('bank-connection-service');

logger.info('Processing sync request', {
  correlationId: 'abc-123',
  service: 'bank-connection-service',
  module: 'F2',
  action: 'sync_transactions'
}, { connectionId: 'conn_123' });
\end{minted}

\subsubsection{Metrics avec Prometheus}

\begin{minted}{typescript}
import { Counter, Histogram, register } from 'prom-client';

// Définition des métriques
export const metrics = {
  httpRequestsTotal: new Counter({
    name: 'http_requests_total',
    help: 'Total number of HTTP requests',
    labelNames: ['method', 'route', 'status_code']
  }),

  httpRequestDuration: new Histogram({
    name: 'http_request_duration_seconds',
    help: 'Duration of HTTP requests in seconds',
    labelNames: ['method', 'route'],
    buckets: [0.1, 0.5, 1, 2, 5, 10]
  }),

  bankSyncOperations: new Counter({
    name: 'bank_sync_operations_total',
    help: 'Total number of bank sync operations',
    labelNames: ['bank_id', 'status']
  }),

  fraudScans: new Counter({
    name: 'fraud_scans_total',
    help: 'Total number of fraud scans',
    labelNames: ['scan_type', 'result']
  })
};

// Middleware de métriques
export function metricsMiddleware(req: any, res: any, next: any) {
  const start = Date.now();
  
  res.on('finish', () => {
    const duration = (Date.now() - start) / 1000;
    
    metrics.httpRequestsTotal.inc({
      method: req.method,
      route: req.route?.path || req.path,
      status_code: res.statusCode
    });
    
    metrics.httpRequestDuration.observe({
      method: req.method,
      route: req.route?.path || req.path
    }, duration);
  });
  
  next();
}
\end{minted}

\begin{bestpractice}
Le monitoring structuré avec des métriques Prometheus et des logs JSON permet une observabilité complète de la plateforme, facilitant le debugging et l'optimisation des performances.
\end{bestpractice}

% ==============================================================================
% CHAPITRE 4 : CONCLUSION
% ==============================================================================
\chapter{Conclusion}

\section{Synèse des Réalisations}

Ce rapport a présenté la conception et l'analyse architecturale d'une Plateforme SaaS d'Intelligence Financière globale, avec un focus particulier sur les modules \module{F2 -- Multi Banking Hub} et \module{F5 -- AI Fraud Detection}.

\subsection{Contributions Techniques}

Mon travail a permis de :

\begin{itemize}
    \item \textbf{Concevoir une architecture modulaire} pour le Multi Banking Hub, supportant les protocoles PSD2 et EBICS via des adaptateurs normalisés.
    \item \textbf{Implémenter une détection de fraude hybride} combinant règles AML/CFT, Machine Learning et analyse de graphes relationnels.
    \item \textbf{Définir des contrats d'interface OpenAPI} robustes avec gestion uniforme des erreurs et support de l'asynchronisme.
    \item \textbf{Appliquer les meilleures pratiques} de développement logiciel : idempotence, sécurité OAuth 2.0/JWT, circuit breakers, et monitoring structuré.
\end{itemize}

\subsection{Intégration dans l'Écosystème}

Les modules F2 et F5 jouent un rôle central dans l'écosystème de la plateforme :

\begin{itemize}
    \item Le \module{F2} alimente l'ensemble des modules en données bancaires centralisées et synchronisées.
    \item Le \module{F5} assure la sécurité transversale en validant les transactions et en détectant les schémas de fraude complexes.
    \item L'intégration avec les modules des collègues (F4, F6, F7) crée une synergie où chaque module bénéficie des données enrichies par les autres.
\end{itemize}

\section{Perspectives d'Évolution}

Plusieurs pistes d'amélioration pourraient être explorées pour renforcer davantage la plateforme :

\begin{itemize}
    \item \textbf{Extension des protocoles bancaires} : Support de protocoles supplémentaires (SWIFT, Host-to-Host).
    \item \textbf{Amélioration du ML} : Intégration de modèles de deep learning pour la détection de fraude.
    \item \textbf{Edge Computing} : Déploiement de composants en edge computing pour réduire la latence.
    \item \textbf{Zero Trust Architecture} : Renforcement de la sécurité avec une approche Zero Trust.
    \item \textbf{Multi-cloud} : Déploiement multi-cloud pour une meilleure résilience.
\end{itemize}

\section{Conclusion Générale}

Ce projet a permis de mettre en œuvre une architecture logicielle moderne et scalable, répondant aux exigences complexes du secteur financier. L'approche modulaire, couplée à des pratiques de développement rigoureuses, assure la maintenabilité et l'évolutivité de la plateforme.

L'intégration harmonieuse des différents modules démontre l'importance d'une vision architecturale globale dans la conception de systèmes SaaS complexes. La plateforme ainsi conçue constitue une base solide pour l'automatisation financière intelligente et la conformité réglementaire.

\begin{center}
\textit{Ce rapport conclut ma mission de développement et d'analyse architecturale au sein de DATAERA,}
\textit{marquant une contribution significative à la modernisation des solutions financières.}
\end{center}

% ==============================================================================
% BIBLIOGRAPHIE
% ==============================================================================
\chapter*{Bibliographie}
\addcontentsline{toc}{chapter}{Bibliographie}

\begin{enumerate}
    \item European Banking Authority. (2019). \textit{Regulatory Technical Standards on Strong Customer Authentication and Common Secure Communication}.
    \item Berlin Group. (2023). \textit{NextGenPSD2 Framework Specification}.
    \item Deutsche Kreditwirtschaft. (2022). \textit{EBICS Specification Version 3.0}.
    \item Financial Action Task Force (FATF). (2023). \textit{International Standards on Combating Money Laundering and the Financing of Terrorism}.
    \item Newman, S. (2021). \textit{Building Microservices: Designing Fine-Grained Systems}. O'Reilly Media.
    \item Richardson, L., & Amundsen, M. (2022). \textit{RESTful Web APIs}. O'Reilly Media.
\end{enumerate}

\end{document}